% આ ભાગનો પૂર્ણ સંપાદનીય પાઠ છે. શૈલી, ફોન્ટ, ચિત્રો અને પ્રકરણસંદર્ભ માટે સંપૂર્ણ સ્રોત ZIP તથા reader/full-reader.tex જુઓ. આ એકલી સ્વતંત્ર કમ્પાઇલ થતી મુખ્ય ફાઇલ નથી.

% BEGIN OLP-0383
\olpart{mvl}{બહુમૂલ્ય તર્કશાસ્ત્ર}
\phantomsection\label{guunit:OLP-0383}


\begin{editorial}
  આ ભાગમાં વિધાનાત્મક બહુમૂલ્ય તર્કશાસ્ત્ર વિશેની
  પ્રારૂપ સામગ્રી છે.
\end{editorial}

\begingroup
\makeatletter\def\ol@sectioncs{section}\makeatother

% BEGIN OLP-0384
\olchapter{mvl}{syn}{વાક્યરચના અને અર્થવિજ્ઞાન}
\olchapterid{mvl}{syn}
\phantomsection\label{guunit:OLP-0384}


\begingroup
\makeatletter\def\ol@sectioncs{section}\makeatother

% BEGIN OLP-0385
\olfileid{mvl}{syn}{int}

\olsection{પરિચય}
\phantomsection\label{guunit:OLP-0385}


શાસ્ત્રીય તર્કશાસ્ત્રમાં આપણે $\lnot$, $\land$, $\lor$,
$\lif$ અને $\liff$ એ વિધાનાત્મક સંયોજકો વડે
પ્રતિજ્ઞાપન ચલોમાંથી બનેલાં સૂત્રોનો
અભ્યાસ કરીએ છીએ. શાસ્ત્રીય તર્કશાસ્ત્રનું અર્થવિજ્ઞાન
વ્યાખ્યાયિત કરતી વખતે $\True$ અને $\False$ એ બે
સત્યમૂલ્યો વાપરીએ છીએ. સત્યમૂલ્ય-નિયુક્તિ~$\pAssign{v}$માં
પ્રતિજ્ઞાપન ચલોને અર્થ આપીએ છીએ; તે
પ્રતિજ્ઞાપન ચલોને $\True$, $\False$ એ
સત્યમૂલ્યો નિયુક્ત કરે છે. દરેક સત્યમૂલ્ય-નિયુક્તિ પછી
કોઈ પણ સૂત્ર~$!A$ માટે સત્યમૂલ્ય
$\pValue{v}(!A)$ નક્કી કરે છે. જ્યારે
$\pValue{v}(!A) = \True$ હોય ત્યારે અને માત્ર
ત્યારે સૂત્ર, સત્યમૂલ્ય-નિયુક્તિ~$\pAssign{v}$માં
સંતોષાય છે; તેને $\pSat{v}{!A}$ લખીએ છીએ.

બહુમૂલ્ય તર્કશાસ્ત્રો $\True$ અને $\False$ ઉપરાંત
વધુ સત્યમૂલ્યો માન્ય રાખીને શાસ્ત્રીય દ્વિમૂલ્ય
તર્કશાસ્ત્રનું સામાન્યકરણ કરે છે. તેથી બહુમૂલ્ય
તર્કશાસ્ત્રમાં સત્યમૂલ્ય-નિયુક્તિ~$\pAssign{v}$ એ
દરેક પ્રતિજ્ઞાપન ચલ~$p$ને શક્ય
સત્યમૂલ્યોમાંથી એક નિયુક્ત કરતું વિધેય છે. માન્ય
સત્યમૂલ્યોના ગણને સામાન્ય રીતે~$V$ કહીશું.
શાસ્ત્રીય તર્કશાસ્ત્ર પણ $V = \{\True, \False\}$
વાળું બહુમૂલ્ય તર્કશાસ્ત્ર છે; તેમાં સંયોજકો માટેની
પરિચિત સત્યતાસારણીઓ વડે સત્યમૂલ્ય
$\pValue{v}(!A)$ ગણાય છે.

વધારાનાં સત્યમૂલ્યો ઉમેરીએ પછી ``અને'', ``અથવા'',
``નહિ'' અને ``જો---તો'' તરીકે વાંચાતા સંયોજકો માટે
$\pValue{v}(!A)$ કેવી રીતે ગણવું તેની એકથી વધુ
સ્વાભાવિક પસંદગીઓ મળે છે. તેથી બહુમૂલ્ય તર્કશાસ્ત્ર
માત્ર સત્યમૂલ્યોના ગણથી જ નહીં, પરંતુ દરેક
સંયોજક માટે પસંદ કરેલાં \emph{સત્યવિધેયો}થી પણ
નક્કી થાય છે. પરંતુ બહુમૂલ્ય તર્કશાસ્ત્ર~$\Log L$
માટે આ પસંદગી થઈ જાય પછી, શાસ્ત્રીય તર્કશાસ્ત્રની
જેમ જ, મૂલ્યાંકન સત્યમૂલ્ય
$\pValue{v}(!A)[\Log L]$ને એકમાત્ર રીતે નક્કી કરે
છે. શાસ્ત્રીય તર્કશાસ્ત્રની જેમ બહુમૂલ્ય તર્કશાસ્ત્રો
પણ \emph{સત્યતાફલનલક્ષી} છે.

અર્થવિજ્ઞાનના આ મૂળ ઘટકો મળ્યા પછી પુનરુક્તિ,
ફલિતતા અને સંતોષ્યતા જેવા અર્થવિજ્ઞાનના ખ્યાલોનો
અનુરૂપ અર્થ વ્યાખ્યાયિત કરી શકીએ. શાસ્ત્રીય
તર્કશાસ્ત્રમાં કોઈ પણ~$\pAssign{v}$ માટે સત્યમૂલ્ય
$\pValue{v}(!A) = \True$ હોય તો સૂત્ર
પુનરુક્તિ છે. બહુમૂલ્ય તર્કશાસ્ત્રમાં આને પણ થોડું
સામાન્ય બનાવવું પડે છે. સૌપ્રથમ, સત્યમૂલ્યોના
ગણ~$V$માં $\True$ હોવું જરૂરી નથી. દાખલા તરીકે,
કેટલાક બહુમૂલ્ય તર્કશાસ્ત્રો સત્યમૂલ્યોના ગણ તરીકે
$0$ અને~$1$ વચ્ચેની બધી પરિમેય સંખ્યાઓ જેવી
સંખ્યાઓ વાપરે છે. એવા કિસ્સામાં $1$ સામાન્ય રીતે
$\True$ની ભૂમિકા ભજવે છે. અન્ય તર્કશાસ્ત્રોમાં
માત્ર એક નહીં, અનેક સત્યમૂલ્યો એવી ભૂમિકા
ભજવે છે. તેથી દરેક બહુમૂલ્ય તર્કશાસ્ત્રમાં
\emph{નિર્દિષ્ટ મૂલ્યો}નો ગણ~$V^+$ હોય એવી શરત
રાખીએ છીએ. પછી $\pValue{v}(!A)[\Log L] \in V^+$
હોય ત્યારે અને માત્ર ત્યારે સૂત્ર,
સત્યમૂલ્ય-નિયુક્તિ~$\pAssign{v}$માં સંતોષાય છે;
તેને $\pSat{v}{!A}[\Log L]$ લખીએ છીએ.
દરેક~$\pAssign{v}$ માટે $\pValue{v}(!A) \in
V^+$ હોય તો અને તો જ સૂત્ર~$!A$ એ
તર્કશાસ્ત્રની પુનરુક્તિ છે, જેને $\Entails[\Log L]
!A$ લખીએ છીએ. અંતે, સૂત્રોના ગણ~$\Gamma$માંથી
$!A$ ફલિત થાય છે એમ કહીએ અને $\Gamma
\Entails[\Log L] !A$ લખીએ છીએ, જો
તે ગણનાં બધાં સૂત્રોને સંતોષતું દરેક
સત્યમૂલ્ય-નિયુક્તિ~$!A$ને પણ સંતોષે.
% END OLP-0385
\endgroup


\begingroup
\makeatletter\def\ol@sectioncs{section}\makeatother

% BEGIN OLP-0386
\olfileid{mvl}{syn}{con}

\olsection{ભાષાઓ અને સંયોજકો}
\phantomsection\label{guunit:OLP-0386}


શાસ્ત્રીય વિધાનાત્મક તર્કશાસ્ત્ર અને અન્ય ઘણાં
તર્કશાસ્ત્રો \emph{વિધાનાત્મક અચલો} તથા
\emph{સંયોજકો}નો ગણ વાપરે છે. દાખલા તરીકે,
નીચેનાંને આપણે મૂળભૂત માનીએ છીએ:
\begin{enumerate}
\item અસત્ય માટેનો વિધાનાત્મક અચલ~$\lfalse$.
\item સત્ય માટેનો વિધાનાત્મક અચલ~$\ltrue$.
\item તાર્કિક સંયોજકો:
  \startycommalist
  \ycomma $\lnot$ (નિષેધ)%
  \ycomma $\land$ (સંયોજન)%
  \ycomma $\lor$ (વિયોજન)%
  \ycomma $\lif$ (પ્રેરણ)%
  \ycomma $\liff$ (દ્વિમુખી પ્રેરણ)%
\end{enumerate}


બહુમૂલ્ય તર્કશાસ્ત્રોમાં પણ આ જ સંયોજકો વપરાય છે.
જોકે, એક જ તર્કશાસ્ત્રમાં, માનો કે સંયોજનનાં જુદાં
સ્વરૂપો સમાવવાં ઘણી વાર ઉપયોગી છે; તેમને અલગ
ઓળખાવવા જુદા સંકેતો જોઈએ. કેટલાક બહુમૂલ્ય
તર્કશાસ્ત્રોમાં શાસ્ત્રીય તર્કશાસ્ત્રમાં સમકક્ષ ન હોય
એવા સંયોજકો પણ છે. તેથી આપણે સામાન્ય કરતાં
થોડી વધુ વ્યાપક રીત લઈશું.

\begin{defn}
  \emph{વિધાનાત્મક ભાષા} એટલે સંયોજકોનો ગણ~$\Lang L$.
  દરેક સંયોજક~$\star$ની \emph{આર્ગ્યુમેન્ટ-સંખ્યા}
  હોય છે; આર્ગ્યુમેન્ટ-સંખ્યા~$n$ ધરાવતો સંયોજક
  \emph{$n$-સ્થાનીય} કહેવાય છે. આર્ગ્યુમેન્ટ-સંખ્યા~$0$
  ધરાવતા સંયોજકો \emph{અચલો} પણ કહેવાય છે;
  આર્ગ્યુમેન્ટ-સંખ્યા~$1$ ધરાવતાં સંયોજકોને
  \emph{એકસ્થાનીય} અને આર્ગ્યુમેન્ટ-સંખ્યા~$2$
  ધરાવતાંને \emph{દ્વિસ્થાનીય} કહે છે.
\end{defn}

\begin{ex}
  વિધાનાત્મક તર્કશાસ્ત્રની પ્રમાણભૂત ભાષા~$\Lang L_0$માં
  નીચેના સંયોજકો છે (કૌંસમાં તેમની આર્ગ્યુમેન્ટ-સંખ્યા છે):
  $\lfalse$~($0$),
  $\lnot$~($1$),
  $\land$~($2$),
  $\lor$~($2$),
  $\lif$~($2$). આપણે વિચારવાના મોટાભાગના
  તર્કશાસ્ત્રો આ ભાષા વાપરશે. કેટલાક તર્કશાસ્ત્રો
  રૂઢિ અનુસાર અમુક સંયોજકો માટે જુદા સંકેતો
  વાપરે છે. દાખલા તરીકે, ગુણન તર્કશાસ્ત્રમાં
  સંયોજનનો સંકેત ઘણી વાર $\land$ને બદલે~$\odot$
  હોય છે. ક્યારેક નવો કારક ઉમેરવો અનુકૂળ હોય છે;
  દાખલા તરીકે, નિર્ધારિતતા કારક~$\triangle$
  ($1$-સ્થાનીય).
\end{ex}
% END OLP-0386
\endgroup


\begingroup
\makeatletter\def\ol@sectioncs{section}\makeatother

% BEGIN OLP-0387
\olfileid{mvl}{syn}{fml}

\olsection{સૂત્રો}
\phantomsection\label{guunit:OLP-0387}


\begin{defn}[સૂત્ર]
\ollabel{defn:formulas}
વિધાનાત્મક ભાષા~$\Lang L$નાં \emph{સૂત્રો}નો
ગણ~$\Frm[L]$ નીચે મુજબ અનુમાનાત્મક રીતે વ્યાખ્યાયિત થાય છે:
\begin{enumerate}
\item દરેક પ્રતિજ્ઞાપન ચલ~$\Obj p_i$ આણ્વિક
  સૂત્ર છે.
\item $\Lang L$નો દરેક $0$-સ્થાનીય સંયોજક
  (વિધાનાત્મક અચલ) આણ્વિક સૂત્ર છે.
\item $\star$ એ $\Lang L$નો $n$-સ્થાનીય સંયોજક
  હોય અને $!A_1$, \dots, $!A_n$ સૂત્રો હોય,
  તો $\star(!A_1, \dots, !A_n)$ પણ સૂત્ર છે.
\item આ સિવાય બીજું કંઈ સૂત્ર નથી.
\end{enumerate}
$\star$ $1$-સ્થાનીય હોય તો $\star(!A_1)$ને ઘણી વાર
માત્ર $\star !A_1$ લખીશું. $\star$ $2$-સ્થાનીય હોય તો
$\star(!A_1,!A_2)$ને ઘણી વાર $(!A_1 \star !A_2)$
લખીશું.
\end{defn}

હંમેશની જેમ, ઘણી વાર સૌથી બહારનાં કૌંસ ચૂપચાપ
છોડી દઈશું.

\begin{ex}
  પ્રમાણભૂત ભાષા~$\Lang{L_0}$માં $\Obj p_1 \lif (\Obj p_1
  \land \lnot \Obj p_2)$ એક સૂત્ર છે. ગુણન
  તર્કશાસ્ત્રની ભાષામાં તેને $\Obj p_1 \lif (\Obj p_1 \odot
  \lnot \Obj p_2)$ તરીકે લખાશે. ભાષામાં $1$-સ્થાનીય
  $\triangle$ ઉમેરીએ તો $\triangle (\Obj p_1
  \land \Obj p_2) \lif (\triangle \Obj p_1 \land \triangle \Obj p_2)$
  જેવાં સૂત્રો પણ મળશે.
\end{ex}
% END OLP-0387
\endgroup


\begingroup
\makeatletter\def\ol@sectioncs{section}\makeatother

% BEGIN OLP-0388
\olfileid{mvl}{syn}{mat}

\olsection{શ્રેણિકો}
\phantomsection\label{guunit:OLP-0388}


બહુમૂલ્ય તર્કશાસ્ત્ર તેની ભાષા, સત્યમૂલ્યોના ગણ~$V$,
નિર્દિષ્ટ સત્યમૂલ્યોના ઉપગણ અને તેના સંયોજકોનાં
સત્યવિધેયો વડે વ્યાખ્યાયિત થાય છે. આ ઘટકોને
એકસાથે \emph{શ્રેણિક} કહે છે.

\begin{defn}[શ્રેણિક]
\ollabel{defn:matrix}
તર્કશાસ્ત્ર~$\Log L$ માટેની \emph{શ્રેણિક}માં આ હોય છે:
\begin{enumerate}
\item ભાષા~$\Lang L$ બનાવતો સંયોજકોનો ગણ;
\item સત્યમૂલ્યોનો અરિક્ત ગણ~$V \neq \emptyset$;
\item નિર્દિષ્ટ સત્યમૂલ્યોનો ગણ~$V^+ \subseteq V$;
\item $\Lang L$ના દરેક $n$-સ્થાનીય સંયોજક~$\star$
માટે સત્યવિધેય~$\tf{\star} : V^n \to V$. જો
$n = 0$ હોય, તો $\tf{\star}$ એ $V$નો માત્ર એક ઘટક છે.
\end{enumerate}
\end{defn}

\begin{ex}
શાસ્ત્રીય તર્કશાસ્ત્ર~\LogCL{} માટેની શ્રેણિકમાં આ હોય છે:
\begin{enumerate}
  \item $\lfalse$, $\lnot$, $\land$, $\lor$ અને $\lif$
  ધરાવતી પ્રમાણભૂત વિધાનાત્મક ભાષા~$\Lang L_0$.
  \item સત્યમૂલ્યોનો ગણ~$V = \{\True, \False\}$.
  \item એકમાત્ર નિર્દિષ્ટ મૂલ્ય~$\True$ છે, એટલે
  $V^+ = \{\True\}$.
  \item $\lfalse$ માટે $\tf{\lfalse} = \False$ છે.
  અન્ય સત્યવિધેયો સામાન્ય સત્યતાસારણીઓ વડે
  અપાય છે (જુઓ \olref{fig:tf-CL}).
\end{enumerate}
\begin{figure}
  \begin{center}
      \begin{tabular}{c|c}
        $\tf{\lnot}$ & \\
        \hline
        $\True$ & $\False$ \\
        $\False$ & $\True$
      \end{tabular}
      \quad
      \begin{tabular}{c|cc}
        $\tf{\land}$ & $\True$ & $\False$ \\
        \hline
        $\True$ & $\True$ & $\False$ \\
        $\False$ & $\False$ & $\False$
      \end{tabular}
      \quad
      \begin{tabular}{c|cc}
        $\tf{\lor}$ & $\True$ & $\False$ \\
        \hline
        $\True$ & $\True$ & $\True$ \\
        $\False$ & $\True$ & $\False$
      \end{tabular}
      \quad
      \begin{tabular}{c|cc}
        $\tf{\lif}$ & $\True$ & $\False$ \\
        \hline
        $\True$ & $\True$ & $\False$ \\
        $\False$ & $\True$ & $\True$
      \end{tabular}
    \end{center}
    \caption{શાસ્ત્રીય તર્કશાસ્ત્ર~$\LogCL$ માટેનાં સત્યવિધેયો.}
    \ollabel{fig:tf-CL}
  \end{figure}
  \end{ex}
% END OLP-0388
\endgroup


\begingroup
\makeatletter\def\ol@sectioncs{section}\makeatother

% BEGIN OLP-0389
\olfileid{mvl}{syn}{val}

\olsection{સત્યમૂલ્ય-નિયુક્તિઓ અને સંતોષ}
\phantomsection\label{guunit:OLP-0389}


\begin{defn}[સત્યમૂલ્ય-નિયુક્તિઓ]
$V$ સત્યમૂલ્યોનો ગણ હોય. $\Lang{L}$ માટે $V$માં મૂલ્યો ધરાવતું
\emph{સત્યમૂલ્ય-નિયુક્તિ} એ ભાષાના પ્રતિજ્ઞાપન ચલોને $V$નો
ઘટક નિયુક્ત કરતું વિધેય~$\pAssign{v}$ છે; એટલે કે,
$\pAssign{v} \colon \PVar \to V$.
\end{defn}

\begin{defn}\ollabel{defn:pValue}
બહુમૂલ્ય તર્કશાસ્ત્ર~$\Log L$ના સત્યમૂલ્યોના ગણ~$V$માં મૂલ્યો
ધરાવતું સત્યમૂલ્ય-નિયુક્તિ~$\pAssign{v}$ આપેલું હોય, ત્યારે મૂલ્યાંકન
વિધેય $\pValue{v} \colon \Frm[L] \to V$ને નીચે પ્રમાણે અનુમાનાત્મક
રીતે વ્યાખ્યાયિત કરો:
  \begin{enumerate}
    \item $\pValue{v}(\Obj p_n) = \pAssign{v}(\Obj p_n)$;
    \item જો $\star$ એ $0$-સ્થાનીય સંયોજક હોય, તો $\pValue{v}(\star)  = \tf{\star}[\Log L]$;
    \item જો $\star$ એ $n$-સ્થાનીય સંયોજક હોય, તો
    \[
      \pValue{v}(\star(!A_1, \dots, !A_n)) = \tf{\star}[\Log L]
      (\pValue{v}(!A_1), \dots, \pValue{v}(!A_n)).
    \]
  \end{enumerate}
\end{defn}

\begin{defn}[સંતોષ]
\ollabel{defn:satisfaction} સૂત્ર~$!A$નો સત્યમૂલ્ય-નિયુક્તિ~$\pAssign{v}$
વડે \emph{સંતોષ} થાય છે, $\pSat{v}{!A}[\Log L]$, ત્યારે જ જ્યારે
$\pValue{v}(!A)[\Log L] \in V^+$ હોય; અહીં $V^+$ એ~$\Log L$ના
નિર્દિષ્ટ સત્યમૂલ્યોનો ગણ છે.

``$\pSat{v}{!A}[\Log L]$ નથી'' એવો અર્થ દર્શાવવા આપણે
$\pSat/{v}{!A}[\Log L]$ લખીએ છીએ. જો $\Gamma$ સૂત્રોનો ગણ હોય, તો
દરેક~$!A \in \Gamma$ માટે $\pSat{v}{!A}[\Log L]$ હોય ત્યારે જ
$\pSat{v}{\Gamma}[\Log L]$.
\end{defn}
% END OLP-0389
\endgroup


\begingroup
\makeatletter\def\ol@sectioncs{section}\makeatother

% BEGIN OLP-0390
\olfileid{mvl}{syn}{sem}

\olsection{અર્થવિચારી ખ્યાલો}
\phantomsection\label{guunit:OLP-0390}


ધારો કે બહુમૂલ્ય તર્કશાસ્ત્ર~$\Log L$ શ્રેણિક વડે આપેલું છે. ત્યારે
આપણે~$\Log L$ માટેના સામાન્ય અર્થવિચારી ખ્યાલો વ્યાખ્યાયિત કરી શકીએ.

\begin{defn}
\begin{enumerate}
\item કોઈ~$\pAssign{v}$ માટે $\pSat{v}{!A}$ હોય, તો
  સૂત્ર~$!A$ \emph{સંતોષ્ય} છે; કોઈ પણ $\pAssign{v}$ માટે
  $\pSat{v}{!A}$ ન હોય, તો તે \emph{અસંતોષ્ય} છે;
\item બધાં સત્યમૂલ્ય-નિયુક્તિઓ~$v$ માટે $\pSat{v}{!A}$ હોય, તો
  સૂત્ર~$!A$ \emph{પુનરુક્તિ} છે;
\item જો $\Gamma$ એ સૂત્રોનો ગણ હોય, તો $\Gamma \Entails !A$
  (``$\Gamma$, $!A$ને ફલિત કરે છે'') ત્યારે અને માત્ર ત્યારે, જ્યારે
  $\pSat{v}{\Gamma}$ હોય એવા દરેક સત્યમૂલ્ય-નિયુક્તિ~$\pAssign{v}$ માટે
  $\pSat{v}{!A}$ હોય.
\item જો $\Gamma$ એ સૂત્રોનો ગણ હોય, તો એવા કોઈ
  સત્યમૂલ્ય-નિયુક્તિ~$\pAssign{v}$ માટે $\pSat{v}{\Gamma}$ હોય ત્યારે
  $\Gamma$ \emph{સંતોષ્ય} છે; અન્યથા $\Gamma$ \emph{અસંતોષ્ય} છે.
\end{enumerate}
\end{defn}

આ ખ્યાલો માટે શાસ્ત્રીય તર્કશાસ્ત્રના કિસ્સા જેવી કેટલીક હકીકતો મળે છે:

\begin{prop}
\ollabel{prop:semanticalfacts}
\begin{enumerate}
\item $!A$ પુનરુક્તિ હોય ત્યારે અને માત્ર ત્યારે
  $\emptyset \Entails !A$;
\item જો $\Gamma$ સંતોષ્ય હોય, તો $\Gamma$નો દરેક સાન્ત ઉપગણ પણ
  સંતોષ્ય છે;
\item\ollabel{def:monotonicity}%
એકદિશવર્ધિતા: જો $\Gamma \subseteq \Delta$
  અને $\Gamma \Entails !A$, તો $\Delta \Entails !A$ પણ;
\item\ollabel{def:Cut}%
પરંપરિતતા: જો $\Gamma \Entails !A$ અને
  $\Delta \cup \{ !A\} \Entails !B$, તો $\Gamma \cup \Delta \Entails
  !B$;
\end{enumerate}
\end{prop}

\begin{proof}
અભ્યાસપ્રશ્ન.
\end{proof}

\begin{prob}
\olref[mvl][syn][sem]{prop:semanticalfacts} સાબિત કરો.
\end{prob}

શાસ્ત્રીય તર્કશાસ્ત્રમાં ફલિતતા અને પ્રેરણ સંયોજક વચ્ચે સંબંધ સ્થાપી શકાય છે.
દાખલા તરીકે, \emph{મોદુસ પોનેન્સ}નું પ્રામાણ્ય મળે છે: જો $\Gamma
\Entails !A$ અને $\Gamma \Entails !A \lif !B$, તો $\Gamma \Entails
!B$. શાસ્ત્રીય તર્કશાસ્ત્રમાં $\Entails$ અને $\lif$ વચ્ચેનો બીજો
મહત્ત્વનો સંબંધ અર્થાનુસારી નિગમન પ્રમેય છે: $\Gamma \Entails !A
\lif !B$ ત્યારે અને માત્ર ત્યારે, જ્યારે $\Gamma \cup \{!A\} \Entails
!B$. આ પરિણામો બહુમૂલ્ય તર્કશાસ્ત્રોમાં \emph{હંમેશાં લાગુ પડતાં નથી}.
તેઓ લાગુ પડે છે કે કેમ તે સત્યવિધેય~$\tf{\lif}$ પર આધાર રાખે છે.
% END OLP-0390
\endgroup


\begingroup
\makeatletter\def\ol@sectioncs{section}\makeatother

% BEGIN OLP-0391
\olfileid{mvl}{syn}{sub}

\olsection{શાસ્ત્રીય તર્કશાસ્ત્ર~$\LogCL$ના ઉપતર્કશાસ્ત્રો તરીકે બહુમૂલ્ય તર્કશાસ્ત્રો}
\phantomsection\label{guunit:OLP-0391}


સામાન્ય બહુમૂલ્ય તર્કશાસ્ત્રોની વ્યાખ્યા એવી શ્રેણિકો વડે થાય છે જેમાં
$\{\True, \False\}$માંથી લેવાયેલા આર્ગ્યુમેન્ટ માટે સત્યવિધેયનું
મૂલ્ય શાસ્ત્રીય સત્યવિધેયના મૂલ્ય સાથે મળે છે. ચોક્કસ કહીએ તો, આ
તર્કશાસ્ત્રોમાં $x \in \{\True, \False\}$ હોય ત્યારે
$\tf{\lnot}[\Log L](x) = \tf{\lnot}[\LogCL](x)$; અને $\star$ એ
$\land$, $\lor$, $\lif$માંથી કોઈ પણ એક હોય અને
$x, y \in \{\True, \False\}$ હોય ત્યારે
$\tf{\star}[\Log L](x,y) = \tf{\star}[\LogCL](x,y)$.
બીજા શબ્દોમાં, $\lnot$, $\land$, $\lor$, $\lif$ માટેનાં સત્યવિધેયો
$\{\True,\False\}$ પર મર્યાદિત કરીએ તો તે ચોક્કસ શાસ્ત્રીય
સત્યવિધેયો છે.

\begin{prop}\ollabel{prop:mvl-cl}
  ધારો કે બહુમૂલ્ય તર્કશાસ્ત્ર~$\Log L$ની ભાષામાં $\lnot$, $\land$,
  $\lor$, $\lif$ સંયોજકો છે, $\True, \False \in V$ છે, અને તેનાં
  સત્યવિધેયો નીચેની શરતો સંતોષે છે:
  \begin{enumerate}
    \item\ollabel{prop:not} $\tf{\lnot}[\Log L](x) = \tf{\lnot}[\LogCL](x)$ જો $x =
    \True$ અથવા $x = \False$;
    \item\ollabel{prop:land} $\tf{\land}[\Log L](x,y) = \tf{\land}[\LogCL](x,y)$,
    \item\ollabel{prop:lor} $\tf{\lor}[\Log L](x,y) = \tf{\lor}[\LogCL](x,y)$,
    \item\ollabel{prop:lif} $\tf{\lif}[\Log L](x,y) = \tf{\lif}[\LogCL](x,y)$,
    જો $x, y \in \{\True, \False\}$ હોય.
  \end{enumerate}
  ત્યારે ઉપરના ચાર સંયોજકો અને વિધાનચલોથી બનેલા કોઈ પણ સૂત્ર માટે,
  તે સૂત્રના દરેક વિધાનચલને શાસ્ત્રીય મૂલ્ય આપતું~$V$માંનું
  કોઈ પણ મૂલ્યાંકન $\pAssign v$, જેના માટે
  $\pAssign v(p) \in \{\True,\False\}$ હોય, તે હેઠળ
  $\pValue v[\Log L](!A) = \pValue v[\LogCL](!A)$.
\end{prop}
\sourcecorrection{OLMV-001}{સ્થિર અંગ્રેજી મૂળ આ સમતા
બહુમૂલ્ય ભાષાના દરેક સૂત્ર માટે કહે છે, છતાં વધારાના
સંયોજકોનાં સત્યવિધેયો પર કોઈ શરત મૂકતું નથી. તે
મૂલ્યાંકનની દ્વિમૂલ્ય શરત પણ એક મુક્ત વિધાનચલ માટે
જ લખે છે. અહીં વિધાનને દર્શાવેલા ચાર સંયોજકોના
સામાન્ય ખંડ અને સૂત્રના દરેક વિધાનચલ સુધી મર્યાદિત
કર્યું છે; શૂન્યસ્થાનીય સંયોજકનો સમાવેશ કરવા તેના
મૂલ્યની અલગ સમાનતા-શરત જોઈએ. દાખલા તરીકે, બીજા
તર્કમાં અસત્ય-અચલને સત્ય મૂલ્ય આપીએ તો ચાર
સંયોજકોની શરતો છતાં આ સમતા તૂટે છે.}

\begin{proof}
  $!A$ પર અનુમાનથી.
  \begin{enumerate}
  \item જો આણ્વિક સૂત્ર $!A \ident p$ હોય, તો
  $\pValue v[\Log L](!A) = \pAssign v(p) = \pValue
  v[\LogCL](!A)$.
  \item જો $!A \ident \lnot B$ હોય, તો
  \begin{align*}
    \pValue v[\Log L](!A) & = \tf{\lnot}[\Log L](\pValue v[\Log L](!B)) & &\text{\olref[val]{defn:pValue} અનુસાર}\\
    & = \tf{\lnot}[\Log L](\pValue v[\LogCL](!B)) && \text{અનુમાનની પૂર્વધારણાથી}\\
    & = \tf{\lnot}[\LogCL](\pValue v[\LogCL](!B))
&&\text{ધારણા \olref{prop:not} અનુસાર,}\\
&&&\text{કારણ કે $\pValue v[\LogCL](!B) \in \{\True, \False\}$,}\\
& = \pValue v[\LogCL](!A) &&\text{\olref[val]{defn:pValue} અનુસાર}.
\end{align*}
\item જો $!A \ident (!B \land !C)$ હોય, તો
\begin{align*}
  \pValue v[\Log L](!A) & = \tf{\land}[\Log L](\pValue v[\Log L](!B), \pValue v[\Log L](!C)) & &\text{\olref[val]{defn:pValue} અનુસાર}\\
  & = \tf{\land}[\Log L](\pValue v[\LogCL](!B),\pValue v[\LogCL](!C)) && \text{અનુમાનની પૂર્વધારણાથી}\\
  & = \tf{\land}[\LogCL](\pValue v[\LogCL](!B),\pValue v[\LogCL](!C))
&&\text{ધારણા \olref{prop:land} અનુસાર,}\\
&&&\text{કારણ કે $\pValue v[\LogCL](!B),\pValue v[\LogCL](!C) \in \{\True, \False\}$,}\\
& = \pValue v[\LogCL](!A) &&\text{\olref[val]{defn:pValue} અનુસાર}.
\end{align*}
\end{enumerate}
$!A \ident (!B \lor !C)$ અને $!A \ident (!B \lif !C)$ના
કિસ્સા સમાન રીતે સાબિત થાય છે.
\end{proof}

\begin{cor}
  જો બહુમૂલ્ય તર્કશાસ્ત્ર \olref{prop:mvl-cl}ની શરતો સંતોષે,
  $\True \in V^+$ અને $\False \notin V^+$ હોય, તો ઉપરના
  સામાન્ય સૂત્ર-ખંડ પર ${\Entails[\Log L]} \subseteq {\Entails[\LogCL]}$;
  એટલે કે, જો $\Gamma \Entails[\Log L] !B$ હોય, તો
  $\Gamma \Entails[\LogCL] !B$ પણ. ખાસ કરીને,
  $\Log L$ની આ ખંડમાં આવેલી દરેક પુનરુક્તિ શાસ્ત્રીય પુનરુક્તિ પણ છે.
\end{cor}
\sourcecorrection{OLMV-002}{સ્થિર અંગ્રેજી મૂળનો
ફલિતતા-સમાવેશ ભાષાના બધા સૂત્રો માટે જણાવાયો છે.
અગાઉના વિધાનથી તેનો પુરાવો માત્ર બંને ભાષાના
સામાન્ય ચાર-સંયોજક ખંડમાં જ મળે છે; અહીં
નિષ્કર્ષનો વ્યાપ એ પ્રમાણે સ્પષ્ટ કર્યો છે.}

\begin{proof}
  આપણે વિરુદ્ધવિધાન સાબિત કરીએ. ધારો કે $\Gamma \Entails/[\LogCL] !B$.
  ત્યારે એવું કોઈ સત્યમૂલ્ય-નિયુક્તિ~$\pAssign v\colon \PVar \to
  \{\True, \False\}$ છે કે દરેક $!A \in \Gamma$ માટે
  $\pValue v[\LogCL](!A) = \True$ અને
  $\pValue v[\LogCL](!B) = \False$. કારણ કે $\True,
  \False \in V$, એ સત્યમૂલ્ય-નિયુક્તિ~$\pAssign v$ એ~$\Log L$ માટેનું
  સત્યમૂલ્ય-નિયુક્તિ પણ છે. \olref{prop:mvl-cl} અનુસાર, દરેક
  $!A \in \Gamma$ માટે $\pValue v[\Log L](!A) =
  \True$ અને $\pValue v[\Log L](!B) = \False$.
  $\True \in V^+$ અને $\False \notin V^+$ હોવાથી
  $\pAssign v \Entails[\Log L] \Gamma$ અને
  $\pAssign v \Entails/[\Log L] !B$; એટલે કે,
  $\Gamma \Entails/[\Log L] !B$.
\end{proof}
% END OLP-0391
\endgroup


\OLEndChapterHook
% END OLP-0384
\endgroup


\begingroup
\makeatletter\def\ol@sectioncs{section}\makeatother

% BEGIN OLP-0392
\olchapter{mvl}{thr}{ત્રિમૂલ્ય તર્કશાસ્ત્રો}
\olchapterid{mvl}{thr}
\phantomsection\label{guunit:OLP-0392}


\begingroup
\makeatletter\def\ol@sectioncs{section}\makeatother

% BEGIN OLP-0393
\olfileid{mvl}{thr}{int}

\olsection{પરિચય}
\phantomsection\label{guunit:OLP-0393}


$\True$ અને $\False$માં માત્ર એક વધુ મૂલ્ય~$\Undef$ ઉમેરીએ,
તો ત્રિમૂલ્ય તર્કશાસ્ત્ર મળે છે. એક જ વધારાનું સત્યમૂલ્ય
હોવા છતાં $\lnot$, $\land$, $\lor$ અને $\lif$ માટે
સત્યવિધેયો વ્યાખ્યાયિત કરવાની ઘણી રીતો છે. પછી કોઈ
તર્કશાસ્ત્ર આ સત્યવિધેયોના કોઈ પણ સંયોજનનો ઉપયોગ કરી શકે છે;
ઉપરાંત માત્ર $\True$ને નિર્દિષ્ટ ગણવું કે $\True$ અને~$\Undef$
બંનેને નિર્દિષ્ટ ગણવાં, એ પસંદગી પણ છે.

અહીં આપણે સૌથી જાણીતા ત્રિમૂલ્ય તર્કશાસ્ત્રોમાંથી કેટલાક,
તેમની પાછળની પ્રેરણા અને તેમના કેટલાક ગુણધર્મો રજૂ કરીએ છીએ.
% END OLP-0393
\endgroup


\begingroup
\makeatletter\def\ol@sectioncs{section}\makeatother

% BEGIN OLP-0394
\olfileid{mvl}{thr}{luk}

\olsection{લુકાશેવિચનું તર્કશાસ્ત્ર}
\phantomsection\label{guunit:OLP-0394}


બહુમૂલ્ય તર્કશાસ્ત્રની પ્રકાશિત અને વિગતે વિકસાવેલી શરૂઆતની
દરખાસ્તોમાંની એક પોલૅન્ડના તત્ત્વચિંતક યાન~\L ukasiewiczએ
1921માં આપી હતી. લુકાશેવિચને પ્રેરણા એરિસ્ટોટલની દરિયાઈ યુદ્ધની
સમસ્યામાંથી મળી હતી: \emph{આજે} ``આવતી કાલે દરિયાઈ યુદ્ધ થશે'' એ
વિધાન ન તો સાચું લાગે છે, ન તો ખોટું; તેનું સત્યમૂલ્ય હજી નક્કી થયું
નથી. લુકાશેવિચે આવાં ``ભાવિ અનિશ્ચિત'' વિધાનો માટે ત્રીજું સત્યમૂલ્ય
રજૂ કરવાની દરખાસ્ત કરી.
\begin{quote}
આવતા વર્ષે કોઈ નિશ્ચિત ક્ષણે, દાખલા તરીકે 21 ડિસેમ્બરે બપોરે,
હું વૉર્સોમાં હાજર હોઈશ કે નહિ, તે અત્યારે હકારાત્મક કે નકારાત્મક
રીતે નક્કી નથી—આ હું વિરોધાભાસ વિના માની શકું છું. તેથી તે
સમયે હું વૉર્સોમાં હોઉં એ શક્ય છે, પણ આવશ્યક નથી. આ ધારણા હેઠળ
``આવતા વર્ષે 21 ડિસેમ્બરે બપોરે હું વૉર્સોમાં હોઈશ'' એ વિધાન અત્યારે
સાચું પણ નથી અને ખોટું પણ નથી. જો તે અત્યારે સાચું હોય, તો મારી
ભાવિ હાજરી આવશ્યક બને, જે ધારણાથી વિરુદ્ધ છે. બીજી બાજુ, જો તે
અત્યારે ખોટું હોય, તો મારી ભાવિ હાજરી અશક્ય બને, જે પણ ધારણાથી
વિરુદ્ધ છે. તેથી વિચારાધીન વિધાન અત્યારે સાચું કે ખોટું નથી; તેને
``0'' અથવા અસત્ય અને ``1'' અથવા સત્યથી જુદું ત્રીજું મૂલ્ય હોવું
જોઈએ. આ મૂલ્યને આપણે ``$\frac{1}{2}$'' વડે દર્શાવી શકીએ. તે
``શક્ય''ને દર્શાવે છે અને ``સાચું'' તથા ``ખોટું'' સાથે ત્રીજા મૂલ્ય
તરીકે જોડાય છે.
\end{quote}
લુકાશેવિચના ત્રીજા સત્યમૂલ્ય માટે આપણે $\Undef$ વાપરીશું.\footnote{લુકાશેવિચ
અહીં ``શક્ય''નો આજકાલ ઓછો પ્રચલિત અર્થ લે છે: શક્ય, પણ આવશ્યક નહિ.}

$\lnot$, $\land$ અને $\lor$ સંયોજકોનાં સત્યવિધેયો આ અર્થઘટન હેઠળ
નક્કી કરવાં સહેલાં છે: ભાવિ અનિશ્ચિત વિધાનનો નિષેધ પણ ભાવિ અનિશ્ચિત
વિધાન છે, તેથી $\tf{\lnot}(\Undef) = \Undef$. સંયોજનનો એક ઘટક
અનિર્ધારિત અને બીજો સાચો હોય, તો આખું સંયોજન પણ અનિર્ધારિત છે—કારણ
કે ભાવિ અનિશ્ચિત ઘટક છેવટે કેવો નીકળે તેના આધારે સંયોજન સાચું પણ
નીકળી શકે અને ખોટું પણ. તેથી \[
    \tf{\land}(\True, \Undef) =
\tf{\land}(\Undef, \True) =
\Undef.
\]
પરંતુ બીજો ઘટક ખોટો હોય તો સંયોજન સાચું નીકળી શકે નહિ, તેથી
\[\tf{\land}(\False, \Undef) =
\tf{\land}(\Undef, \False) = \False.\]
\sourcecorrection{OLMV-003}{સ્થિર અંગ્રેજી મૂળના આ
સમીકરણમાં બંને ડાબી બાજુએ એક જ ક્રમમાં અસત્ય અને
અનિર્ધારિત મૂલ્ય લખાયું છે. નીચેની સંયોજન-સારણી
અનુસાર બીજો ક્રમ ઉલટો છે; અહીં બંને ક્રમ દર્શાવ્યા છે.}
બાકીનાં મૂલ્યો માટે, જ્યારે આર્ગ્યુમેન્ટનાં સત્યમૂલ્યો $\True$
અથવા $\False$ તરીકે નક્કી હોય, ત્યારે શાસ્ત્રીય તર્કશાસ્ત્રના
જ મૂલ્યો મળે છે.

પ્રેરણ સંયોજકનો કિસ્સો થોડો વધુ મુશ્કેલ છે. ધારો કે $q$ ભાવિ
અનિશ્ચિત વિધાન છે. જો $p$ ખોટું હોય, તો $q$ છેવટે ગમે તે નીકળે,
$p \lif q$ સાચું રહેશે; તેથી $\tf{\lif}(\False, \Undef) = \True$
રાખવું જોઈએ. અને જો $p$ સાચું હોય, તો $q$ ગમે તે નીકળે,
$q \lif p$ સાચું રહેશે; તેથી $\tf{\lif}(\Undef, \True) = \True$.
જો $p$ સાચું હોય, તો $p \lif q$ છેવટે સાચું કે ખોટું નીકળી શકે;
તેથી $\tf{\lif}(\True, \Undef) = \Undef$. એ જ રીતે, જો $p$
ખોટું હોય, તો $q \lif p$ છેવટે સાચું કે ખોટું નીકળી શકે; તેથી
$\tf{\lif}(\Undef, \False) = \Undef$. હવે $p$ અને $q$ બંને
ભાવિ અનિશ્ચિત હોય એવો કિસ્સો બાકી રહે છે. પ્રેરણાના આધાર પરથી
આ કિસ્સામાં ખરેખર $\Undef$ આપવું જોઈએ. પરંતુ એમ કરવાથી
$!A \lif !A$ \emph{પુનરુક્તિ નહિ રહે}. લુકાશેવિચને
$!A \lor \lnot !A$ અને $\lnot(!A \land \lnot !A)$ છોડવામાં
અડચણ નહોતી, પણ $!A \lif !A$ છોડવામાં વાંધો હતો. તેથી તેણે
$\tf{\lif}(\Undef, \Undef) = \True$ નક્કી કર્યું.

\begin{defn}\ollabel{def:lukasiewicz}
ત્રિમૂલ્ય લુકાશેવિચ તર્કશાસ્ત્ર નીચેની શ્રેણિક વડે વ્યાખ્યાયિત થાય છે:
\begin{enumerate}
  \item $\lnot$, $\land$, $\lor$, $\lif$ ધરાવતી પ્રમાણભૂત
  વિધાનાત્મક ભાષા~$\Lang L_0$.
  \item સત્યમૂલ્યોનો ગણ $V = \{\True, \Undef, \False\}$.
  \item એકમાત્ર નિર્દિષ્ટ મૂલ્ય $\True$ છે; એટલે $V^+ = \{\True\}$.
  \item અસત્ય-અચલનું મૂલ્ય અસત્ય છે. અન્ય સત્યવિધેયો નીચેની
  સારણીઓ વડે અપાય છે:
  \begin{center}
    \begin{tabular}{c|c}
      $\tf{\lnot}$ & \\
      \hline
      $\True$ & $\False$ \\
      $\Undef$ & $\Undef$ \\
      $\False$ & $\True$
    \end{tabular}
    \quad
    \begin{tabular}{c|ccc}
      $\tf{\land}[\LogLuk[3]]$ & $\True$ & $\Undef$ & $\False$ \\
      \hline
      $\True$ & $\True$ & $\Undef$ & $\False$ \\
      $\Undef$ & $\Undef$ & $\Undef$ & $\False$\\
      $\False$ & $\False$ & $\False$ & $\False$
    \end{tabular}
    \\[2ex]
    \begin{tabular}{c|ccc}
      $\tf{\lor}[\LogLuk[3]]$ & $\True$ & $\Undef$ & $\False$ \\
      \hline
      $\True$ & $\True$ & $\True$ & $\True$ \\
      $\Undef$ & $\True$ & $\Undef$ & $\Undef$ \\
      $\False$ & $\True$ & $\Undef$ & $\False$
    \end{tabular}
    \quad
    \begin{tabular}{c|ccc}
      $\tf{\lif}[\LogLuk[3]]$ & $\True$ & $\Undef$ & $\False$ \\
      \hline
      $\True$ & $\True$ & $\Undef$ & $\False$ \\
      $\Undef$ & $\True$ & $\True$ & $\Undef$ \\
      $\False$ & $\True$ & $\True$ & $\True$
    \end{tabular}
  \end{center}
\end{enumerate}
\end{defn}
\sourcecorrection{OLMV-004}{પહેલાં $\Lang L_0$માં
અસત્ય-અચલ શૂન્યસ્થાનીય સંયોજક તરીકે વ્યાખ્યાયિત છે,
પણ સ્થિર અંગ્રેજી મૂળની આ શ્રેણિક તેનો સત્યવિધેય
આપતી નથી. શાસ્ત્રીય અર્થ જાળવવા અહીં તેનું
મૂલ્ય અસત્ય હોવાનું સ્પષ્ટ કર્યું છે.}

સરળતાથી જોઈ શકાય છે કે માત્ર $\lnot$, $\land$ અને $\lor$
ધરાવતું કોઈ પણ સૂત્ર~$!A$, જો તેનાં બધાં
પ્રતિજ્ઞાપન ચલોને $\Undef$ આપવામાં આવે, તો તેનું
સત્યમૂલ્ય~$\Undef$ થશે. તેથી, દાખલા તરીકે, શાસ્ત્રીય પુનરુક્તિઓ
$p \lor \lnot p$ અને $\lnot(p \land \lnot p)$ એ
$\LogLuk[3]$માં પુનરુક્તિ નથી, કારણ કે $\pAssign v(p) = \Undef$
હોય ત્યારે $\pValue{v}(!A) = \Undef$.

જ્યાં $\pAssign v(p) = \True$ અથવા $\False$ હોય એવા
સત્યમૂલ્ય-નિયુક્તિઓ પર $\pValue v(!A)$ તેનું શાસ્ત્રીય સત્યમૂલ્ય જ આપશે.

\begin{prop}
  $!A$માં આવતા દરેક $p$ માટે જો $\pAssign v(p) \in \{\True, \False\}$,
  તો $\pValue v(!A)[\LogLuk[3]] = \pValue v(!A)[\LogCL]$.
\end{prop}

\begin{prob}\label{mvl:thr:luk:prob:luk-iff} ધારો કે~$\LogLuk[3]$માં
  $\pValue v(!A \liff !B) = \pValue v((!A \lif !B) \land (!B \lif
  !A))$ એમ વ્યાખ્યાયિત કરીએ. ત્યારે $\liff$ની સત્યતાસારણી કેવી હશે?
\end{prob}

ઘણી શાસ્ત્રીય પુનરુક્તિઓ \LogLuk[3]માં પણ પુનરુક્તિ \emph{છે};
દાખલા તરીકે $\lnot p \lif (p \lif q)$. શાસ્ત્રીય તર્કશાસ્ત્રની
જેમ સત્યતાસારણીથી આ તપાસી શકીએ છીએ:
\begin{center}
\begin{tabular}{cc|cccccc}
  $p$ & $q$ & $\lnot$ & $p$ & $\lif$ & $\smash{(}p$ & $\lif$ & $q\smash{)}$ \\ \hline
  \True & \True & \False & \True & \True & \True & \True & \True \\
  \True & \Undef & \False & \True & \True & \True & \Undef & \Undef \\
  \True & \False & \False & \True & \True & \True & \False & \False \\
  \Undef & \True & \Undef & \Undef & \True & \Undef & \True & \True \\
  \Undef & \Undef & \Undef & \Undef & \True & \Undef & \True & \Undef \\
  \Undef & \False & \Undef & \Undef & \True & \Undef & \Undef & \False \\
  \False & \True & \True & \False & \True & \False & \True & \True \\
  \False & \Undef & \True & \False & \True & \False & \True & \Undef \\
  \False & \False & \True & \False & \True & \False & \True & \False \\
\end{tabular}
\end{center}

\begin{prob}
  નીચેનાં સૂત્રો \LogLuk[3]માં પુનરુક્તિ છે તે બતાવો:
  \begin{enumerate}
    \item $p \lif (q \lif p)$
    \item\label{mvl:thr:luk:prob:luk-taut-2} $\lnot(p \land q) \liff (\lnot p \lor \lnot q)$
    \item\label{mvl:thr:luk:prob:luk-taut-3} $\lnot(p \lor q) \liff (\lnot p \land \lnot q)$
  \end{enumerate}
  (\olref[mvl][thr][luk]{prob:luk-taut-2} અને
  \olref[mvl][thr][luk]{prob:luk-taut-3}માં $!A \liff !B$ને
  $(!A \lif !B) \land (!B \lif !A)$નું સંક્ષિપ્ત લેખન માનો,
  અથવા \olref[mvl][thr][luk]{prob:luk-iff}ના તમારા ઉકેલનો ઉપયોગ કરો.)
\end{prob}

\begin{prob}
  નીચેનાં શાસ્ત્રીય પુનરુક્તિઓ \LogLuk[3]માં પુનરુક્તિ નથી તે બતાવો:
  \begin{enumerate}
    \item $(\lnot p \land p) \lif q$
    \item $((p \lif q) \lif p) \lif p$
    \item $(p \lif (p \lif q)) \lif (p \lif q)$
  \end{enumerate}
\end{prob}
\sourcecorrection{OLMV-005}{સ્થિર અંગ્રેજી મૂળના
પહેલા અભ્યાસપ્રશ્નમાં અંતે વધારાનો જમણો કૌંસ છે.
અહીં વધારાનો કૌંસ દૂર કર્યો છે; સૂત્રનો અર્થ અને
શાસ્ત્રીય પુનરુક્તિરૂપ જળવાય છે.}

આથી કદાચ એમ લાગી શકે કે બધી શાસ્ત્રીય પુનરુક્તિઓ
$\LogLuk[3]$માં પુનરુક્તિ ન હોય તોપણ દરેક
સત્યમૂલ્ય-નિયુક્તિ હેઠળ તેમનું મૂલ્ય ઓછામાં ઓછું $\True$ અથવા
$\Undef$ હશે. પરંતુ એવું નથી. નીચેનું પ્રતિઉદાહરણ જુઓ:
\[
  \lnot(p \lif \lnot p) \lor \lnot(\lnot p \lif p)
\]
જો $p$નું મૂલ્ય~$\Undef$ હોય, તો આ સૂત્રનું મૂલ્ય $\False$ છે.

\begin{prob}
  લુકાશેવિચ તર્કશાસ્ત્રમાં નીચેનામાંથી કયા ફલિતતા-સંબંધો સાચા છે?
  દરેક માટે સત્યતાસારણી આપો.
  \begin{enumerate}
    \item $p, p \lif q \Entails q$
    \item $\lnot\lnot p \Entails p$
    \item $p \land q \Entails p$
    \item $p \Entails p \land p$
    \item $p \Entails p \lor q$
  \end{enumerate}
\end{prob}

લુકાશેવિચે પોતાના ત્રિમૂલ્ય તર્કશાસ્ત્રના આધારે સંભાવનાનું
તર્કશાસ્ત્ર બનાવવાની આશા રાખી હતી. તેણે એકસ્થાનીય
સંયોજક $\Diamond !A$ (``$!A$ શક્ય છે'' માટે) અને તેને અનુરૂપ
$\Box !A$ (``$!A$ આવશ્યક છે'' માટે) રજૂ કર્યા:
\begin{center}
  \begin{tabular}{c|c}
    $\tf{\Diamond}$ & \\
    \hline
    $\True$ & $\True$ \\
    $\Undef$ & $\True$ \\
    $\False$ & $\False$
  \end{tabular}
  \quad
  \begin{tabular}{c|c}
    $\tf{\Box}$ & \\
    \hline
    $\True$ & $\True$ \\
    $\Undef$ & $\False$ \\
    $\False$ & $\False$
  \end{tabular}
\end{center}
બીજા શબ્દોમાં, $p$ શક્ય ત્યારે અને માત્ર ત્યારે છે જ્યારે
તે પહેલેથી ખોટું હોવાનું નક્કી નથી; અને $p$ આવશ્યક ત્યારે
અને માત્ર ત્યારે છે જ્યારે તે પહેલેથી સાચું હોવાનું નક્કી છે.

\begin{prob}
  $\Box p \liff \lnot\Diamond \lnot p$ અને $\Diamond p \liff
  \lnot \Box \lnot p$ એ $\Box$ અને~$\Diamond$ની ઉપરની
  સત્યતાસારણીઓથી વિસ્તૃત \LogLuk[3]માં પુનરુક્તિ છે તે બતાવો.
\end{prob}

પરંતુ સૂચિત આ વિધાત્મક તર્કશાસ્ત્રની મર્યાદા ટૂંક સમયમાં
સ્પષ્ટ થઈ: ભવિષ્યમાં ગમે તે થાય, $p \land \lnot p$ કદી સાચું
નીકળી શકે નહિ. તેથી અત્યારે તેનું મૂલ્ય અનિર્ધારિત હોય તોપણ તેને
અશક્ય ગણવું જોઈએ; એટલે $\lnot \Diamond(p \land \lnot p)$ પુનરુક્તિ
હોવું જોઈએ. પરંતુ જો $\pAssign v(p) = \Undef$, તો
$\pValue v(\lnot \Diamond(p \land \lnot p)) = \False$.
\sourcecorrection{OLMV-006}{સ્થિર અંગ્રેજી મૂળ અહીં
અંતિમ મૂલ્ય અનિર્ધારિત કહે છે. ઉપરની સારણીઓ મુજબ
અનિર્ધારિત વિધાન અને તેનો નિષેધ ભેગાં કરીએ તો
અનિર્ધારિત મળે છે, તેનું સંભાવના-મૂલ્ય સત્ય અને
તેનો નિષેધ અસત્ય મળે છે. તેથી અહીં અસત્ય લખ્યું છે;
પ્રતિઉદાહરણનો હેતુ યથાવત્ રહે છે.}
લુકાશેવિચે શક્યતા અને આવશ્યકતા જેવા વિધાત્મક ભેદો માટે
બે સત્યમૂલ્યો પૂરતાં નથી એમ સાચું કહ્યું હતું, પરંતુ
ત્રીજું સત્યમૂલ્ય ઉમેરવું પણ પૂરતું નથી.
% END OLP-0394
\endgroup


\begingroup
\makeatletter\def\ol@sectioncs{section}\makeatother

% BEGIN OLP-0395
\olfileid{mvl}{thr}{skl}

\olsection{ક્લીની તર્કશાસ્ત્રો}
\phantomsection\label{guunit:OLP-0395}


સ્ટીફન ક્લીનીએ બે ત્રિમૂલ્ય તર્કશાસ્ત્રો રજૂ કર્યાં. તેમની પ્રેરણા
એવા તર્કશાસ્ત્રમાંથી આવી જેમાં સત્યમૂલ્યોને સંગણનાત્મક પ્રક્રિયાઓના
પરિણામો માનવામાં આવે છે: કોઈ પ્રક્રિયા $\True$ અથવા $\False$ આપી
શકે, પણ તે અટકે જ નહિ એવું પણ બને. ત્યારે તેને અનુરૂપ સત્યમૂલ્ય
અવ્યાખ્યાયિત હોય છે અને તેને~$\Undef$ વડે દર્શાવીએ છીએ.

વિધાન~$!A$નો નિષેધ ગણવા માટે પહેલાં~$!A$નું મૂલ્ય ગણવું પડે અને
પછી તેનું વિપરીત મૂલ્ય આપવું પડે. જો $!A$ની ગણતરી અટકે નહિ, તો
આખી પ્રક્રિયા પણ અટકતી નથી; તેથી $\Undef$નો નિષેધ $\Undef$ છે.

સંયોજન $!A \land !B$ ગણવાની બે રીતો છે. પહેલાં~$!A$ અને પછી
$!B$ ગણીએ, તો બંનેનાં પરિણામ~$\True$ હોય ત્યારે જવાબ $\True$,
અન્યથા $\False$ મળે. બેમાંથી એક પણ ગણતરી અટકે નહિ, તો આખી
પ્રક્રિયા પણ અટકતી નથી. તેથી આ પદ્ધતિમાં એક ઘટક અવ્યાખ્યાયિત
હોય, તો સંયોજન પણ અવ્યાખ્યાયિત છે. વિયોજન માટે પણ એવું જ છે.

પરંતુ $!A$ અને $!B$ને સમાંતરે ગણાવી શકીએ તો વધુ સારું થાય.
બેમાંથી એક પ્રક્રિયા અટકી $\False$ આપે, તો આપણે અટકી શકીએ,
કારણ કે જવાબ ખોટો જ હશે. એટલે આ કિસ્સામાં એક ઘટક ખોટો હોય
તો બીજો અવ્યાખ્યાયિત હોવા છતાં સંયોજન ખોટું છે. એવી જ રીતે,
વિયોજન સમાંતરે ગણતાં કોઈ એક ઘટક સાચો નીકળે એટલે અટકી શકાય:
વિયોજન સાચું જ હશે. આમ એક ઘટક અવ્યાખ્યાયિત હોય તોપણ સંયુક્ત
વિધાનનું પરિણામ જાણી શકાય છે. આ અર્થઘટનમાં $\Undef$ને
``અવ્યાખ્યાયિત'' કરતાં ``અજ્ઞાત'' તરીકે વાંચી શકીએ.

આ બે અર્થઘટનોમાંથી ક્લીનીનાં પ્રબળ અને દુર્બળ તર્કશાસ્ત્રો મળે છે.
પ્રેરણ સંયોજકને $\lnot !A \lor !B$ની સમકક્ષ તરીકે વ્યાખ્યાયિત કરીએ છીએ.

\begin{defn}
\emph{પ્રબળ ક્લીની તર્કશાસ્ત્ર}~$\LogKs$ નીચેની શ્રેણિક વડે વ્યાખ્યાયિત થાય છે:
\begin{enumerate}
  \item $\lnot$, $\land$, $\lor$, $\lif$ ધરાવતી અને
  અસત્ય-અચલ વિનાની પ્રમાણભૂત વિધાનાત્મક ભાષા~$\Lang L_0$ની ઉપભાષા.
  \item સત્યમૂલ્યોનો ગણ $V = \{\True, \Undef, \False\}$.
  \item એકમાત્ર નિર્દિષ્ટ મૂલ્ય $\True$ છે; એટલે $V^+ = \{\True\}$.
  \item સત્યવિધેયો નીચેની સારણીઓ વડે અપાય છે:
  \begin{center}
    \begin{tabular}{c|c}
      $\tf{\lnot}$ & \\
      \hline
      $\True$ & $\False$ \\
      $\Undef$ & $\Undef$ \\
      $\False$ & $\True$
    \end{tabular}
    \quad
    \begin{tabular}{c|ccc}
      $\tf{\land}[\LogKs]$ & $\True$ & $\Undef$ & $\False$ \\
      \hline
      $\True$ & $\True$ & $\Undef$ & $\False$ \\
      $\Undef$ & $\Undef$ & $\Undef$ & $\False$\\
      $\False$ & $\False$ & $\False$ & $\False$
    \end{tabular}
    \\[2ex]
    \begin{tabular}{c|ccc}
      $\tf{\lor}[\LogKs]$ & $\True$ & $\Undef$ & $\False$ \\
      \hline
      $\True$ & $\True$ & $\True$ & $\True$ \\
      $\Undef$ & $\True$ & $\Undef$ & $\Undef$ \\
      $\False$ & $\True$ & $\Undef$ & $\False$
    \end{tabular}
    \quad
    \begin{tabular}{c|ccc}
      $\tf{\lif}[\LogKs]$ & $\True$ & $\Undef$ & $\False$ \\
      \hline
      $\True$ & $\True$ & $\Undef$ & $\False$ \\
      $\Undef$ & $\True$ & $\Undef$ & $\Undef$ \\
      $\False$ & $\True$ & $\True$ & $\True$
    \end{tabular}
  \end{center}
\end{enumerate}
\end{defn}

\begin{defn}
\emph{દુર્બળ ક્લીની તર્કશાસ્ત્ર}~$\LogKw$ નીચેની શ્રેણિક વડે વ્યાખ્યાયિત થાય છે:
\begin{enumerate}
  \item $\lnot$, $\land$, $\lor$, $\lif$ ધરાવતી અને
  અસત્ય-અચલ વિનાની પ્રમાણભૂત વિધાનાત્મક ભાષા~$\Lang L_0$ની ઉપભાષા.
  \item સત્યમૂલ્યોનો ગણ $V = \{\True, \Undef, \False\}$.
  \item એકમાત્ર નિર્દિષ્ટ મૂલ્ય $\True$ છે; એટલે $V^+ = \{\True\}$.
  \item સત્યવિધેયો નીચેની સારણીઓ વડે અપાય છે:
  \begin{center}
    \begin{tabular}{c|c}
      $\tf{\lnot}$ & \\
      \hline
      $\True$ & $\False$ \\
      $\Undef$ & $\Undef$ \\
      $\False$ & $\True$
    \end{tabular}
    \quad
    \begin{tabular}{c|ccc}
      $\tf{\land}[\LogKw]$ & $\True$ & $\Undef$ & $\False$ \\
      \hline
      $\True$ & $\True$ & $\Undef$ & $\False$ \\
      $\Undef$ & $\Undef$ & $\Undef$ & $\Undef$\\
      $\False$ & $\False$ & $\Undef$ & $\False$
    \end{tabular}
    \\[2ex]
    \begin{tabular}{c|ccc}
      $\tf{\lor}[\LogKw]$ & $\True$ & $\Undef$ & $\False$ \\
      \hline
      $\True$ & $\True$ & $\Undef$ & $\True$ \\
      $\Undef$ & $\Undef$ & $\Undef$ & $\Undef$ \\
      $\False$ & $\True$ & $\Undef$ & $\False$
    \end{tabular}
    \quad
    \begin{tabular}{c|ccc}
      $\tf{\lif}[\LogKw]$ & $\True$ & $\Undef$ & $\False$ \\
      \hline
      $\True$ & $\True$ & $\Undef$ & $\False$ \\
      $\Undef$ & $\Undef$ & $\Undef$ & $\Undef$ \\
      $\False$ & $\True$ & $\Undef$ & $\True$
    \end{tabular}
  \end{center}
\end{enumerate}
\end{defn}
\sourcecorrection{OLMV-007}{સ્થિર અંગ્રેજી મૂળ બંને
શ્રેણિકોમાં આખી પ્રમાણભૂત ભાષા~$\Lang L_0$ કહે છે;
પરંતુ અગાઉ તેની અંદર અસત્ય-અચલ પણ છે, જ્યારે અહીં
તેનું સત્યવિધેય અપાયું નથી. તેને અસત્ય મૂલ્ય આપીએ
તોપણ તેના નિષેધથી પુનરુક્તિ મળે છે અને નીચેનું
પ્રમેય ખોટું પડે છે. તેથી બંને શ્રેણિકોને અહીં
ચોક્કસ ચાર-સંયોજક, અચલ વિનાની ઉપભાષા સુધી
મર્યાદિત કરી છે.}

\begin{prop}
  $\LogKs$ અને $\LogKw$માં કોઈ પુનરુક્તિ નથી.
\end{prop}

\begin{proof}
  બધાં પ્રતિજ્ઞાપન ચલો~$p$ માટે જો
  $\pAssign v(p) = \Undef$, તો કોઈ પણ સૂત્ર~$!A$નું સત્યમૂલ્ય
  $\pValue v(!A) = \Undef$ હશે, કારણ કે બંને તર્કશાસ્ત્રોમાં
  \[
    \tf{\lnot}(\Undef) = \tf{\lor}(\Undef, \Undef) = \tf{\land}(\Undef,
  \Undef) = \tf{\lif}(\Undef, \Undef) = \Undef
  \]
  અને $\LogKs$ કે $\LogKw$માંથી કોઈ પણ તર્કશાસ્ત્રમાં
  $\Undef \notin V^+$.
  તેથી આ સત્યમૂલ્ય-નિયુક્તિ હેઠળ $!A$ નિર્દિષ્ટ નથી.
\end{proof}

પ્રબળ અને દુર્બળ ક્લીની તર્કશાસ્ત્રોમાં પુનરુક્તિ ન હોવા છતાં
તેમના ફલિતતા-સંબંધો ગૌણ નથી.

\begin{prob}
  નીચેનામાંથી કયા સંબંધો (a)~પ્રબળ અને (b)~દુર્બળ ક્લીની
  તર્કશાસ્ત્રમાં સાચા છે? દરેક માટે સત્યતાસારણી આપો.
  \begin{enumerate}
    \item $p, p \lif q \Entails q$
    \item $p \lor q, \lnot p \Entails q$
    \item $p \land q \Entails p$
    \item $p \Entails p \land p$
    \item $p \Entails p \lor q$
  \end{enumerate}
\end{prob}

દિમિત્રી બોચ્વારે $\Undef$નો અર્થ ``નિરર્થક'' લીધો અને
જૂઠ્ઠાણાના વિરોધાભાસ જેવા વિરોધાભાસો ઉકેલવા માટે વિરોધાભાસી
વિધાનોને~$\Undef$ મૂલ્ય આપવાની દરખાસ્ત કરી. તેણે એવો તર્કશાસ્ત્ર
રજૂ કર્યો જે મૂળભૂત રીતે વધારાના સંયોજકો સાથેનું દુર્બળ ક્લીની
તર્કશાસ્ત્ર છે. એમાંથી બે ``બાહ્ય નિષેધ'' અને ``અવ્યાખ્યાયિત છે''
કારક છે:
\begin{center}
  \begin{tabular}{c|c}
    $\tf{\sim}$ & \\
    \hline
    $\True$ & $\False$ \\
    $\Undef$ & $\True$ \\
    $\False$ & $\True$
  \end{tabular}
\quad
  \begin{tabular}{c|c}
    $\tf{+}$ & \\
    \hline
    $\True$ & $\False$ \\
    $\Undef$ & $\True$ \\
    $\False$ & $\False$
  \end{tabular}
\end{center}

\begin{prob}
  બોચ્વારના તર્કશાસ્ત્રમાં $\lnot$ અને~$+$ વડે $\sim$ વ્યાખ્યાયિત
  કરી શકાય છે? એટલે માત્ર પ્રતિજ્ઞાપન ચલ~$p$ ધરાવતું અને $\sim$ વિનાનું એવું સૂત્ર
  શોધો જેનું સત્યમૂલ્ય હંમેશાં~$\mathord{\sim}p$ જેટલું જ હોય.
  તમારો જવાબ સાચો છે તે દર્શાવવા સત્યતાસારણી આપો.
\end{prob}
% END OLP-0395
\endgroup


\begingroup
\makeatletter\def\ol@sectioncs{section}\makeatother

% BEGIN OLP-0396
\olfileid{mvl}{thr}{god}

\olsection{ગોડેલ તર્કશાસ્ત્રો}
\phantomsection\label{guunit:OLP-0396}


કુર્ત ગોડેલે $n$-મૂલ્ય તર્કશાસ્ત્રોની એવી શ્રેણી રજૂ કરી કે
તેમાંના દરેકમાં સ્ફુરણાવાદી તર્કશાસ્ત્રમાં પ્રામાણ્ય ધરાવતાં બધાં
સૂત્રો સામેલ છે, અને તે દરેક તર્કશાસ્ત્ર શાસ્ત્રીય
તર્કશાસ્ત્રમાં સમાયેલું છે.
તેમાંનું પહેલું રસપ્રદ તર્કશાસ્ત્ર આ છે:

\begin{defn}\ollabel{defn:goedel}
\emph{$3$-મૂલ્ય ગોડેલ તર્કશાસ્ત્ર}~$\LogGod$ નીચેની શ્રેણિક વડે
વ્યાખ્યાયિત થાય છે:
\begin{enumerate}
  \item $\lfalse$, $\lnot$, $\land$, $\lor$, $\lif$
  ધરાવતી પ્રમાણભૂત વિધાનાત્મક ભાષા~$\Lang L_0$.
  \item સત્યમૂલ્યોનો ગણ $V = \{\True, \Undef, \False\}$.
  \item એકમાત્ર નિર્દિષ્ટ મૂલ્ય $\True$ છે; એટલે $V^+ = \{\True\}$.
  \item $\lfalse$ માટે $\tf{\lfalse} = \False$ છે. અન્ય
  સંયોજકોનાં સત્યવિધેયો નીચેની સારણીઓ વડે અપાય છે:
  \begin{center}
    \begin{tabular}{c|c}
      $\tf{\lnot}[\LogGod]$ & \\
      \hline
      $\True$ & $\False$ \\
      $\Undef$ & $\False$ \\
      $\False$ & $\True$
    \end{tabular}
    \quad
    \begin{tabular}{c|ccc}
      $\tf{\land}[\LogGod]$ & $\True$ & $\Undef$ & $\False$ \\
      \hline
      $\True$ & $\True$ & $\Undef$ & $\False$ \\
      $\Undef$ & $\Undef$ & $\Undef$ & $\False$\\
      $\False$ & $\False$ & $\False$ & $\False$
    \end{tabular}
    \\[2ex]
    \begin{tabular}{c|ccc}
      $\tf{\lor}[\LogGod]$ & $\True$ & $\Undef$ & $\False$ \\
      \hline
      $\True$ & $\True$ & $\True$ & $\True$ \\
      $\Undef$ & $\True$ & $\Undef$ & $\Undef$ \\
      $\False$ & $\True$ & $\Undef$ & $\False$
    \end{tabular}
    \quad
    \begin{tabular}{c|ccc}
      $\tf{\lif}[\LogGod]$ & $\True$ & $\Undef$ & $\False$ \\
      \hline
      $\True$ & $\True$ & $\Undef$ & $\False$ \\
      $\Undef$ & $\True$ & $\True$ & $\False$ \\
      $\False$ & $\True$ & $\True$ & $\True$
    \end{tabular}
  \end{center}
\end{enumerate}
\end{defn}

અહીં $\land$ અને~$\lor$ની સત્યતાસારણીઓ લુકાશેવિચ અને પ્રબળ
ક્લીની તર્કશાસ્ત્ર જેવી જ છે, પણ $\lnot$ અને~$\lif$ની
સત્યતાસારણીઓ દરેકમાં જુદી છે. ગોડેલ તર્કશાસ્ત્રમાં
$\tf{\lnot}(\Undef) = \False$. લુકાશેવિચ અને ક્લીની
તર્કશાસ્ત્રોથી વિપરીત, $\tf{\lif}(\Undef, \False) = \False$;
અને ક્લીની તર્કશાસ્ત્રથી વિપરીત, પરંતુ લુકાશેવિચ તર્કશાસ્ત્રની જેમ,
$\tf{\lif}(\Undef, \Undef) = \True$.

ઉપર સૂચવેલા સ્ફુરણાવાદી તર્કશાસ્ત્ર સાથેના સંબંધ પરથી સમજાય છે કે
$\LogGod[3]$ તેની નજીક છે. બધી સ્ફુરણાવાદી સત્યતાઓ
$\LogGod[3]$માં પુનરુક્તિ છે; જ્યારે સ્ફુરણાવાદી અર્થમાં પ્રામાણ્ય
ન ધરાવતી ઘણી શાસ્ત્રીય પુનરુક્તિઓ $\LogGod[3]$માં પણ પુનરુક્તિ નથી.
દાખલા તરીકે, આ પુનરુક્તિઓ નથી:
\begin{align*}
  & p \lor \lnot p && (p \lif q) \lif (\lnot p \lor q) \\
  & \lnot\lnot p \lif p && \lnot(\lnot p \land \lnot q) \lif (p \lor q) \\
  & ((p \lif q) \lif p) \lif p && \lnot(p \lif q) \lif (p \land \lnot q)
\end{align*}
પરંતુ $\LogGod[3]$ની દરેક પુનરુક્તિ સ્ફુરણાવાદી અર્થમાં પણ
પ્રામાણ્ય ધરાવે એવું નથી; દાખલા તરીકે $\lnot\lnot p \lor \lnot p$
અથવા $(p \lif q) \lor (q \lif p)$.

\begin{prob}
  નીચેનાં સૂત્રો~$\LogGod[3]$માં પુનરુક્તિ છે તે બતાવવા
  સત્યતાસારણીઓ આપો:
  \begin{align*}
    & \lnot\lnot p \lor \lnot p\\
    & (p \lif q) \lor (q \lif p) \\
    & \lnot(p \land q) \lif (\lnot p \lor \lnot q) \\
    & (p \lif q) \lor (q \lif r) \lor (r \lif s)
  \end{align*}
\end{prob}

\begin{prob}
  નીચેનાં સૂત્રો~$\LogGod[3]$માં પુનરુક્તિ નથી તે બતાવવા
  સત્યતાસારણીઓ આપો:
  \begin{align*}
    & (p \lif q) \lif (\lnot p \lor q) \\
    & \lnot(\lnot p \land \lnot q) \lif (p \lor q) \\
    & ((p \lif q) \lif p) \lif p \\
    & \lnot(p \lif q) \lif (p \land \lnot q)
  \end{align*}
\end{prob}

\begin{prob}
  ગોડેલ તર્કશાસ્ત્રમાં નીચેનામાંથી કયા ફલિતતા-સંબંધો સાચા છે?
  દરેક માટે સત્યતાસારણી આપો.
  \begin{enumerate}
    \item $p, p \lif q \Entails q$
    \item $p \lor q, \lnot p \Entails q$
    \item $p \land q \Entails p$
    \item $p \Entails p \land p$
    \item $p \Entails p \lor q$
  \end{enumerate}
\end{prob}
% END OLP-0396
\endgroup


\begingroup
\makeatletter\def\ol@sectioncs{section}\makeatother

% BEGIN OLP-0397
\olfileid{mvl}{thr}{mul}

\olsection{માત્ર $\True$ જ નિર્દિષ્ટ નહિ}
\phantomsection\label{guunit:OLP-0397}


અત્યાર સુધી જોયેલા બધાં તર્કશાસ્ત્રોમાં નિર્દિષ્ટ સત્યમૂલ્યોનો ગણ
$V^+ = \{\True\}$ હતો; એટલે કોઈ વસ્તુનું સત્યમૂલ્ય~$\True$
હોય ત્યારે અને માત્ર ત્યારે જ તેને સાચી ગણાતી. પરંતુ કોઈ વસ્તુ
માત્ર~$\False$ ન હોય તોપણ તેને સાચી ગણી શકાય. દાખલા તરીકે,
શ્રેણિકમાં $V^+ = \{\True, \Undef\}$ નક્કી કરવાથી આવું
તર્કશાસ્ત્ર મળશે.

\begin{defn}
\emph{વિરોધાભાસનું તર્કશાસ્ત્ર}~$\LogLP$ નીચેની શ્રેણિક વડે વ્યાખ્યાયિત થાય છે:
\begin{enumerate}
  \item $\lnot$, $\land$, $\lor$, $\lif$ ધરાવતી અને
  અસત્ય-અચલ વિનાની પ્રમાણભૂત વિધાનાત્મક ભાષા~$\Lang L_0$ની ઉપભાષા.
  \item સત્યમૂલ્યોનો ગણ $V = \{\True, \Undef, \False\}$.
  \item $\True$ અને $\Undef$ નિર્દિષ્ટ છે; એટલે $V^+ = \{\True, \Undef\}$.
  \item સત્યવિધેયો પ્રબળ ક્લીની તર્કશાસ્ત્ર જેવા જ છે.
\end{enumerate}
\end{defn}

\begin{defn}
હાલ્ડેનનું \emph{નિરર્થકતાનું તર્કશાસ્ત્ર}~$\LogHal$ નીચેની શ્રેણિક વડે
વ્યાખ્યાયિત થાય છે:
\begin{enumerate}
  \item $\lnot$, $\land$, $\lor$, $\lif$ અને $1$-સ્થાનીય
  સંયોજક~$+$ ધરાવતી, અસત્ય-અચલ વિનાની પ્રમાણભૂત
  વિધાનાત્મક ભાષા~$\Lang L_0$ની વિસ્તૃત ઉપભાષા.
  \item સત્યમૂલ્યોનો ગણ $V = \{\True, \Undef, \False\}$.
  \item $\True$ અને $\Undef$ નિર્દિષ્ટ છે; એટલે $V^+ = \{\True, \Undef\}$.
  \item સત્યવિધેયો દુર્બળ ક્લીની તર્કશાસ્ત્ર જેવા છે, અને ઉપરાંત
  ``નિરર્થક છે'' કારક આ પ્રમાણે છે:
  \begin{center}
    \begin{tabular}{c|c}
    $\tf{+}$ & \\
    \hline
    $\True$ & $\False$ \\
    $\Undef$ & $\True$ \\
    $\False$ & $\False$
    \end{tabular}
  \end{center}
\end{enumerate}
\end{defn}
\sourcecorrection{OLMV-008}{સ્થિર અંગ્રેજી મૂળ બંને
શ્રેણિકોમાં આખી ભાષા~$\Lang L_0$ કહે છે, પરંતુ
તેમાં અગાઉથી અસત્ય-અચલ છે અને અહીં તેનું સત્યવિધેય
આપેલું નથી. અગાઉના ક્લીની વિભાગની જેમ અહીં
બંને શ્રેણિકોને અચલ વિનાની સામાન્ય ઉપભાષા પર
વ્યાખ્યાયિત કરી છે. વધારાનો કારક ધરાવતા હાલ્ડેન
તર્કશાસ્ત્રની પુનરુક્તિઓની શાસ્ત્રીય સાથેની સરખામણી
પણ સામાન્ય ચાર-સંયોજક ખંડ પૂરતી છે.}

સમાન સત્યતાસારણીઓ ધરાવતા ક્લીની તર્કશાસ્ત્રોથી વિપરીત,
આ બંનેમાં પુનરુક્તિઓ \emph{છે}.

\begin{prop}\ollabel{prop:LP-taut-CL} સામાન્ય ચાર-સંયોજક ખંડમાં
  $\LogLP$ની પુનરુક્તિઓ શાસ્ત્રીય વિધાનાત્મક તર્કશાસ્ત્રની
  પુનરુક્તિઓ જેવી જ છે.
\end{prop}

\begin{proof}
  \olref[syn][sub]{prop:mvl-cl} અનુસાર, જો $\Entails[\LogLP] !A$
  હોય, તો $\Entails[\LogCL] !A$. વિપરીત દિશા માટે બતાવીએ કે
  જો એવું કોઈ સત્યમૂલ્ય-નિયુક્તિ~$\pAssign v\colon \PVar \to \{\False, \True,
  \Undef\}$ હોય કે $\pValue v(!A)[\LogKs] = \False$, તો એવું
  સત્યમૂલ્ય-નિયુક્તિ~$\pAssign {v'}\colon \PVar \to \{\False, \True\}$
  હોય કે $\pValue {v'}(!A)[\LogCL] = \False$. આથી $\LogLP$ માટે
  પરિણામ મળે છે, કારણ કે $\LogKs$ અને $\LogLP$નાં લાક્ષણિક
  સત્યવિધેયો સરખાં છે અને $\LogLP$નું એકમાત્ર અનિર્દિષ્ટ
  સત્યમૂલ્ય $\False$ છે (આ જ $\LogLP$ અને~$\LogKs$ વચ્ચેનો ભેદ છે).
  તેથી જો $\Entails/[\LogLP] !A$, તો કોઈ સત્યમૂલ્ય-નિયુક્તિ~$\pAssign v$
  માટે $\pValue v(!A)[\LogLP] = \pValue
  v[\LogKs](!A) = \False$. આપણે જે દાવો સાબિત કરીએ છીએ તે મુજબ
  $\pValue{v'}[\LogCL](!A) = \False$; એટલે
  $\Entails/[\LogCL] !A$.

  દાવો સાબિત કરવા પહેલાં $\pAssign {v'}$ને આ રીતે વ્યાખ્યાયિત કરીએ:
  \[
    \pAssign {v'}(p) =
    \begin{cases}
    \True & \text{જો } \pAssign {v}(p) \in \{\True, \Undef\}\\
    \False & \text{અન્યથા}
  \end{cases}
  \]
  હવે $!A$ પર અનુમાનથી બતાવીએ કે (a)~જો $\pValue v(!A)[\LogKs]
  = \False$, તો $\pValue {v'}(!A)[\LogCL] = \False$; અને (b)~જો
  $\pValue v(!A)[\LogKs] = \True$, તો $\pValue {v'}(!A)[\LogCL] =
  \True$.
  \begin{enumerate}
    \item અનુમાનનો આધાર: $!A \ident p$. 
    \olref[syn][val]{defn:pValue} અનુસાર,
    $\pValue v(!A)[\LogKs] = \pAssign v(p)$.
    જો આ મૂલ્ય $\False$ અથવા $\True$ હોય, તો ઉપરની
    $\pAssign {v'}$ની વ્યાખ્યા મુજબ $\pValue {v'}(!A)[\LogCL]$
    પણ એ જ મૂલ્ય છે; એટલે (a) અને~(b) બંને મળે છે.
    જો મૂલ્ય $\Undef$ હોય, તો બંને પૂર્વપક્ષ ખોટા છે.
    \sourcecorrection{OLMV-009}{સ્થિર અંગ્રેજી મૂળ
    બધા મૂલ્યાંકનો માટે $\pAssign v(p)=\pValue{v'}(!A)$
    કહે છે; પરંતુ $\pAssign v(p)=\Undef$ હોય ત્યારે
    $\pAssign{v'}(p)=\True$ છે. અહીં માત્ર (a) અને
    (b) માટે જરૂરી બે શરતી દાવા સાબિત કર્યા છે.}

    અનુમાનના પગલા માટે નીચેના કિસ્સા લો:
    \item $!A \ident \lnot !B$.
    \begin{enumerate}
      \item ધારો કે $\pValue v(\lnot!B)[\LogKs] = \False$.
      $\tf{\lnot}[\LogKs]$ની વ્યાખ્યા મુજબ
      $\pValue v(!B)[\LogKs] = \True$. અનુમાનની પૂર્વધારણાના
      કિસ્સા~(b)થી $\pValue {v'}(!B)[\LogCL] = \True$ મળે છે;
      તેથી $\pValue {v'}(\lnot !B)[\LogCL] = \False$.
      \item ધારો કે $\pValue v(\lnot!B)[\LogKs] = \True$.
      $\tf{\lnot}[\LogKs]$ની વ્યાખ્યા મુજબ
      $\pValue v(!B)[\LogKs] = \False$. અનુમાનની પૂર્વધારણાના
      કિસ્સા~(a)થી $\pValue {v'}(!B)[\LogCL] = \False$ મળે છે;
      તેથી $\pValue {v'}(\lnot !B)[\LogCL] = \True$.
    \end{enumerate}

    \item $!A \ident (!B \land !C)$.
    \begin{enumerate}
      \item ધારો કે $\pValue v(!B \land !C)[\LogKs] = \False$.
      $\tf{\land}[\LogKs]$ની વ્યાખ્યા મુજબ
      $\pValue v(!B)[\LogKs] = \False$ અથવા
      $\pValue v(!C)[\LogKs] = \False$. અનુમાનની પૂર્વધારણાના
      કિસ્સા~(a)થી $\pValue {v'}(!B)[\LogCL] = \False$
      અથવા $\pValue {v'}(!C)[\LogCL] = \False$ મળે છે; તેથી
      $\pValue {v'}(!B \land !C)[\LogCL] = \False$.
      \item ધારો કે $\pValue v(!B \land !C)[\LogKs] = \True$.
      $\tf{\land}[\LogKs]$ની વ્યાખ્યા મુજબ
      $\pValue v(!B)[\LogKs] = \True$ અને
      $\pValue v(!C)[\LogKs] = \True$. અનુમાનની પૂર્વધારણાના
      કિસ્સા~(b)થી $\pValue {v'}(!B)[\LogCL] = \True$
      અને $\pValue {v'}(!C)[\LogCL] = \True$ મળે છે; તેથી
      $\pValue {v'}(!B \land !C)[\LogCL] = \True$.
    \end{enumerate}
    \sourcecorrection{OLMV-010}{સ્થિર અંગ્રેજી મૂળ
    સંયોજન ખોટું અને સાચું હોય તે બંને કિસ્સામાં
    બે વાર~$!B$નું મૂલ્ય લખે છે, જોકે આગળનો
    નિષ્કર્ષ $!B$ અને~$!C$ પર આધારિત છે. અહીં
    બીજા ઘટકનું મૂલ્ય $!C$ તરીકે સુધાર્યું છે.}
  \end{enumerate}

  બીજા બે કિસ્સા સમાન રીતે સાબિત થાય છે; તેમને અભ્યાસપ્રશ્ન તરીકે
  રાખ્યા છે. વિકલ્પે, ઉપરનો પુરાવો માત્ર $\lnot$ અને~$\land$
  ધરાવતા બધાં સૂત્રો માટે પરિણામ સ્થાપિત કરે છે. હવે
  $\LogKs$ અને $\LogCL$ બંનેમાં કોઈ પણ $\pAssign v$ માટે
  $\pValue v(!B \lor
  !C) = \pValue v(\lnot(\lnot!B \land \lnot !C))$ અને
  $\pValue v(!B \lif
  !C) = \pValue v(\lnot(!B \land \lnot !C))$ એ હકીકતોનો
  ઉપયોગ કરી શકાય છે.
\end{proof}

\begin{prob}
  \olref[mvl][thr][mul]{prop:LP-taut-CL}નો પુરાવો પૂર્ણ કરો; એટલે
  $!A \ident (!B \lor !C)$ અને $!A \ident (!B \lif !C)$ના
  કિસ્સાઓમાં (a) અને~(b) સ્થાપિત કરો.
\end{prob}

\begin{prob}
દરેક શાસ્ત્રીય પુનરુક્તિ~$\LogHal$માં પુનરુક્તિ છે તે સાબિત કરો.
\end{prob}

સામાન્ય ચાર-સંયોજક ખંડમાં તેમની પુનરુક્તિઓ શાસ્ત્રીય તર્કશાસ્ત્ર
જેવી હોવા છતાં તેમના ફલિતતા-સંબંધો જુદા છે. દાખલા તરીકે,
$\LogLP$ \emph{વિરોધાભાસ-સહનશીલ} છે: $\lnot p, p \Entails/ q$.
તેથી વિસ્ફોટનો સિદ્ધાંત $\lnot !A, !A \Entails !B$ સામાન્ય રીતે
લાગુ પડતો નથી. (કેટલાક $!A$ અને $!B$ના કિસ્સામાં તે લાગુ પડે છે;
દાખલા તરીકે, $!B$~પુનરુક્તિ હોય ત્યારે.)

\begin{prob}
  નીચેનામાંથી કયા સંબંધો (a)~$\LogLP$માં અને (b)~$\LogHal$માં
  સાચા છે? દરેક માટે સત્યતાસારણી આપો.
  \begin{enumerate}
    \item $p, p \lif q \Entails q$
    \item $\lnot q, p \lif q \Entails \lnot p$
    \item $p \lor q, \lnot p \Entails q$
    \item $\lnot p, p \Entails q$
    \item $p \Entails p \lor q$
    \item $p \lif q, q\lif r \Entails p \lif r$
  \end{enumerate}
\end{prob}

$\LogLuk[3]$માં $\Undef$ને નિર્દિષ્ટ ગણીએ તો શું થાય?

\begin{defn}
\emph{ત્રિમૂલ્ય આર-મિંગલ}~$\LogRM[3]$ નીચેની શ્રેણિક વડે વ્યાખ્યાયિત થાય છે:
  \begin{enumerate}
    \item $\lfalse$, $\lnot$, $\land$, $\lor$, $\lif$
    ધરાવતી પ્રમાણભૂત વિધાનાત્મક ભાષા~$\Lang L_0$.
    \item સત્યમૂલ્યોનો ગણ $V = \{\True, \Undef, \False\}$.
    \item $\True$ અને $\Undef$ નિર્દિષ્ટ છે; એટલે $V^+ = \{\True, \Undef\}$.
    \item સત્યવિધેયો લુકાશેવિચ તર્કશાસ્ત્ર~$\LogLuk[3]$ જેવા જ છે.
  \end{enumerate}
\end{defn}

\begin{prob}
  $\LogRM[3]$માં નીચેનામાંથી કયા સંબંધો સાચા છે?
  \begin{enumerate}
    \item $p, p \lif q \Entails q$
    \item $p \lor q, \lnot p \Entails q$
    \item $\lnot p, p \Entails q$
    \item $p \Entails p \lor q$
  \end{enumerate}
\end{prob}

નિર્દિષ્ટ મૂલ્યો બદલવાથી જુદી સત્યતાસારણીઓ ક્યારેક સમાન
તર્કશાસ્ત્ર, એટલે સમાન ફલિતતા-સંબંધ, આપી શકે છે. દાખલા તરીકે,
ગોડેલ તર્કશાસ્ત્રમાં $\{\True\}$ને બદલે
$V^+ = \{\True, \Undef\}$ લઈએ ત્યારે એવું થાય છે.

\begin{prop}\ollabel{prop:gl-udes}
  $V = \{\False, \Undef,\True\}$, $V^+=\{\True,
  \Undef\}$ અને $3$-મૂલ્ય ગોડેલ તર્કશાસ્ત્રનાં સત્યવિધેયો ધરાવતી
  શ્રેણિક શાસ્ત્રીય તર્કશાસ્ત્ર વ્યાખ્યાયિત કરે છે.
\end{prop}

\begin{proof}
  અભ્યાસપ્રશ્ન.
\end{proof}

\begin{prob}
  ગોડેલ તર્કશાસ્ત્રની જેમ વ્યાખ્યાયિત પરંતુ
  $V^+=\{\True,\Undef\}$ ધરાવતા તર્કશાસ્ત્ર~$\Log L$ માટે,
  જો $\Gamma \Entails/[\Log L] !B$ તો
  $\Gamma \Entails/[\LogCL] !B$ છે તે બતાવી
  \olref[mvl][thr][mul]{prop:gl-udes} સાબિત કરો.
  \olref[mvl][thr][mul]{prop:LP-taut-CL}ના વિચારોનો ઉપયોગ કરો;
  પરંતુ (a) અને~(b) સાબિત કરવાને બદલે
  $\pValue v(!A)[\LogGod] = \False$ ત્યારે અને માત્ર ત્યારે જ્યારે
  $\pValue {v'}(!A)[\LogCL] = \False$ છે તે બતાવો
  (અને તેથી $\pValue v(!A)[\LogGod] \in \{\True,\Undef\}$ ત્યારે
  અને માત્ર ત્યારે જ્યારે $\pValue {v'}(!A)[\LogCL] = \True$).
  આથી વિધાન કેવી રીતે સ્થાપિત થાય છે તે સમજાવો.
\end{prob}
% END OLP-0397
\endgroup


\OLEndChapterHook
% END OLP-0392
\endgroup


\begingroup
\makeatletter\def\ol@sectioncs{section}\makeatother

% BEGIN OLP-0398
\olchapter{mvl}{inf}{અનંતમૂલ્ય તર્કશાસ્ત્રો}
\olchapterid{mvl}{inf}
\phantomsection\label{guunit:OLP-0398}


\begingroup
\makeatletter\def\ol@sectioncs{section}\makeatother

% BEGIN OLP-0399
\olfileid{mvl}{inf}{int}

\olsection{પરિચય}
\phantomsection\label{guunit:OLP-0399}


શ્રેણિકમાં સત્યમૂલ્યોની સંખ્યા સાન્ત હોવી જરૂરી નથી. અનંત સત્યમૂલ્યોના
ગણ માટે સહજ પસંદગી $0$ અને~$1$ વચ્ચેની સંમેય સંખ્યાઓનો ગણ
$V_\infty = [0,1] \cap \Rat$ છે; એટલે કે,
\begin{align*}
    V_\infty & = \Setabs{\frac{n}{m}}{n,m \in \Nat \text{ અને } m>0 \text{ તથા } n\le m}.
\intertext{આ અનંત સત્યમૂલ્ય-ગણને ધ્યાનમાં લેતાં, તેના નીચેના ઉપગણો
પણ ધ્યાનમાં લેવા ઘણી વાર ઉપયોગી છે; અહીં $m\ge 2$ છે:}
V_m & = \Setabs{\frac{n}{m-1}}{n \in \Nat \text{ અને } n<m}
\intertext{દાખલા તરીકે, $V_5$ એ સમાન અંતરે આવેલાં $5$ સત્યમૂલ્યોનો ગણ છે,}
V_5 & = \{0, \frac{1}{4}, \frac{1}{2}, \frac{3}{4}, 1\}.
\end{align*}
\sourcecorrection{OLMV-011}{સ્થિર અંગ્રેજી મૂળના
અનંત સંમેય ગણમાં ભાગાકારના છેદ માટે શૂન્યને
બાકાત રાખતું નથી. શૂન્યનો ભાગાકાર અવ્યાખ્યાયિત
હોવાથી અહીં છેદ ધન હોય તે શરત સ્પષ્ટ કરી છે.}
\sourcecorrection{OLMV-012}{સ્થિર અંગ્રેજી મૂળના
સૂત્રમાં $V_m$ માટે અંશને $m$ સુધી મંજૂર કરે છે,
જેથી દર્શાવેલા $V_5$માં પાંચને બદલે છ મૂલ્યો અને
એક કરતાં મોટું મૂલ્ય મળે. અહીં અંશ $m$ કરતાં
નાનો રાખ્યો છે તથા છેદ શૂન્ય ન થાય તે માટે
$m\ge 2$ સ્પષ્ટ કર્યું છે.}
આ સત્યમૂલ્ય-ગણો પર આધારિત તર્કશાસ્ત્રોમાં સામાન્ય રીતે માત્ર
$1$ નિર્દિષ્ટ હોય છે; એટલે $V^+ = \{1\}$. બીજા શબ્દોમાં,
$1$ને (સંપૂર્ણ) સત્ય અને $0$ને સંપૂર્ણ અસત્યની ભૂમિકા આપીએ છીએ,
પરંતુ સૂત્રોનું મૂલ્ય~$V$માંનું કોઈ પણ મધ્યવર્તી મૂલ્ય હોઈ શકે.

આ ઉપરાંત $0$ અને~$1$ વચ્ચેની બધી \emph{વાસ્તવિક} સંખ્યાઓનો ગણ
$V_{[0,1]} = [0,1]$ અથવા $[0,1]$ના બીજા અનંત ઉપગણો પણ વિચારી શકાય.
આ સત્યમૂલ્ય-ગણ ધરાવતા તર્કશાસ્ત્રોને ઘણી વાર \emph{અસ્પષ્ટ}
તર્કશાસ્ત્રો કહે છે.
% END OLP-0399
\endgroup


\begingroup
\makeatletter\def\ol@sectioncs{section}\makeatother

% BEGIN OLP-0400
\olfileid{mvl}{inf}{luk}

\olsection{લુકાશેવિચનું તર્કશાસ્ત્ર}
\phantomsection\label{guunit:OLP-0400}


\begin{editorial}
  અનંતમૂલ્ય લુકાશેવિચ તર્કશાસ્ત્ર અંગેનો આ વિભાગ ટૂંકું
  ``પ્રારૂપ'' માત્ર છે.
\end{editorial}

\begin{defn}\ollabel{def:lukasiewicz} અનંતમૂલ્ય લુકાશેવિચ
તર્કશાસ્ત્ર~$\LogLuk[\infty]$ નીચેની શ્રેણિક વડે વ્યાખ્યાયિત થાય છે:
\begin{enumerate}
  \item $\lnot$, $\land$, $\lor$, $\lif$ ધરાવતી
  પ્રમાણભૂત વિધાનાત્મક ભાષા~$\Lang L_0$.
  \item સત્યમૂલ્યોનો ગણ $V_\infty$.
  \item એકમાત્ર નિર્દિષ્ટ મૂલ્ય $1$ છે; એટલે $V^+ = \{1\}$.
  \item અસત્ય-અચલનું મૂલ્ય શૂન્ય છે. અન્ય સત્યવિધેયો
  નીચેના વિધેયો વડે અપાય છે:
  \begin{align*}
    \tf{\lnot}[\LogLuk](x) & = 1 - x\\
    \tf{\land}[\LogLuk](x,y) & = \min(x,y)\\
    \tf{\lor}[\LogLuk](x,y) & = \max(x,y)\\
    \tf{\lif}[\LogLuk](x,y) & = \min(1,1-(x-y)) = \begin{cases}
      1 & \text{જો } x \le y\\
      1-(x-y) & \text{અન્યથા.}
    \end{cases}
    \end{align*}
\end{enumerate}
$m$-મૂલ્ય લુકાશેવિચ તર્કશાસ્ત્રની વ્યાખ્યા પણ આવી જ છે;
માત્ર $V = V_m$ છે.
\end{defn}
\sourcecorrection{OLMV-013}{સ્થિર અંગ્રેજી મૂળ
અહીં પ્રમાણભૂત ભાષા~$\Lang L_0$ લે છે, જેમાં
અગાઉથી અસત્ય-અચલ સામેલ છે, પરંતુ તેનું સત્યવિધેય
આપતું નથી. અગાઉની ત્રિમૂલ્ય લુકાશેવિચ શ્રેણિક
સાથે એકરૂપતા રાખવા અહીં તેનું મૂલ્ય શૂન્ય
સ્પષ્ટ કર્યું છે.}

\begin{prop}
  \olref[thr][luk]{def:lukasiewicz} વડે વ્યાખ્યાયિત
  $\LogLuk[3]$ એ \olref{def:lukasiewicz} વડે વ્યાખ્યાયિત
  $\LogLuk[3]$ જેવું જ છે.
\end{prop}

\begin{proof}
  \olref[thr][luk]{def:lukasiewicz}માં આપેલી સંયોજકોની
  સત્યતાસારણીઓની સરખામણી \olref{def:lukasiewicz}નાં
  સમીકરણોથી મળતી નીચેની સત્યતાસારણીઓ સાથે કરીને આ જોઈ શકાય છે:
  \begin{center}
    \begin{tabular}{c|c}
      $\tf{\lnot}$ & \\
      \hline
      $1$ & $0$ \\
      $1/2$ & $1/2$ \\
      $0$ & $1$
    \end{tabular}
    \quad
    \begin{tabular}{c|ccc}
      $\tf{\land}[\LogLuk[3]]$ & $1$ & $1/2$ & $0$ \\
      \hline
      $1$ & $1$ & $1/2$ & $0$ \\
      $1/2$ & $1/2$ & $1/2$ & $0$\\
      $0$ & $0$ & $0$ & $0$
    \end{tabular}
    \\[2ex]
    \begin{tabular}{c|ccc}
      $\tf{\lor}[\LogLuk[3]]$ & $1$ & $1/2$ & $0$ \\
      \hline
      $1$ & $1$ & $1$ & $1$ \\
      $1/2$ & $1$ & $1/2$ & $1/2$ \\
      $0$ & $1$ & $1/2$ & $0$
    \end{tabular}
    \quad
    \begin{tabular}{c|ccc}
      $\tf{\lif}[\LogLuk[3]]$ & $1$ & $1/2$ & $0$ \\
      \hline
      $1$ & $1$ & $1/2$ & $0$ \\
      $1/2$ & $1$ & $1$ & $1/2$ \\
      $0$ & $1$ & $1$ & $1$
    \end{tabular}
  \end{center}
\end{proof}

\begin{prop}\ollabel{prop:luk-infty-m}
  જો $\Gamma \Entails[\LogLuk[\infty]] !B$, તો બધાં~$m \ge 2$
  માટે $\Gamma \Entails[\LogLuk[m]] !B$.
\end{prop}

\begin{proof}
  અભ્યાસપ્રશ્ન.
\end{proof}

\begin{prob}
  \olref[mvl][inf][luk]{prop:luk-infty-m} સાબિત કરો.
\end{prob}

હકીકતમાં, સાન્ત આધારગણ માટે વિપરીત દિશા પણ સાચી છે.
\sourcecorrection{OLMV-014}{સ્થિર અંગ્રેજી મૂળ
વિપરીત દિશા કોઈ મર્યાદા વિના કહે છે. સાન્ત
આધારગણ માટે વિરોધી મૂલ્યાંકનમાં આવતા માત્ર
સાન્ત ઘણા સંમેય મૂલ્યોને એક જ સાન્ત જાળીમાં
સમાવી શકાય છે. પરંતુ અનંત આધારગણ જુદા જુદા
છેદવાળાં અસંખ્ય ચલોનાં મૂલ્યો નિર્ધારિત કરી શકે;
બધી સાન્ત જાળીઓમાં તેના આધાર અસંતોષ્ય છતાં
અનંત સંમેય જાળીમાં સંતોષ્ય હોઈ શકે. તેથી અહીં
દાવો સાન્ત આધારગણ સુધી મર્યાદિત કર્યો છે.}

અનંતમૂલ્ય લુકાશેવિચ તર્કશાસ્ત્ર સૌથી પ્રચલિત અસ્પષ્ટ
તર્કશાસ્ત્ર છે. અસ્પષ્ટ તર્કશાસ્ત્રના સાહિત્યમાં પ્રેરણ સંયોજકને
ઘણી વાર $\lnot !A \lor !B$ તરીકે વ્યાખ્યાયિત કરે છે. ત્યારે
અનંતમૂલ્ય પ્રબળ ક્લીની તર્કશાસ્ત્ર મળશે.

\begin{prob}
  $(p \lif q) \lor (q \lif p)$ એ~$\LogLuk[\infty]$ની
  પુનરુક્તિ છે તે બતાવો.
\end{prob}
% END OLP-0400
\endgroup


\begingroup
\makeatletter\def\ol@sectioncs{section}\makeatother

% BEGIN OLP-0401
\olfileid{mvl}{inf}{god}

\olsection{ગોડેલનાં તર્કશાસ્ત્રો}
\phantomsection\label{guunit:OLP-0401}


\begin{editorial}
  અનંતમૂલ્ય ગોડેલ તર્કશાસ્ત્ર અંગેનો આ વિભાગ ટૂંકું
  ``પ્રારૂપ'' માત્ર છે.
\end{editorial}

\begin{defn}\ollabel{def:goedel} અનંતમૂલ્ય ગોડેલ
તર્કશાસ્ત્ર~$\LogGod[\infty]$ નીચેની શ્રેણિક વડે વ્યાખ્યાયિત થાય છે:
\begin{enumerate}
  \item $\lfalse$, $\lnot$, $\land$, $\lor$, $\lif$ ધરાવતી
  પ્રમાણભૂત વિધાનાત્મક ભાષા~$\Lang L_0$.
  \item સત્યમૂલ્યોનો ગણ $V_\infty$.
  \item એકમાત્ર નિર્દિષ્ટ મૂલ્ય $1$ છે; એટલે $V^+ = \{1\}$.
  \item સત્યવિધેયો નીચેના વિધેયો વડે અપાય છે:
  \begin{align*}
    \tf{\lfalse} & = 0\\
    \tf{\lnot}[\LogGod](x) & = \begin{cases}
      1 & \text{જો } x =0\\
      0 & \text{અન્યથા}
    \end{cases}\\
    \tf{\land}[\LogGod](x,y) & = \min(x,y)\\
    \tf{\lor}[\LogGod](x,y) & = \max(x,y)\\
    \tf{\lif}[\LogGod](x,y) & = \begin{cases}
      1 & \text{જો } x \le y\\
      y & \text{અન્યથા.}
    \end{cases}
    \end{align*}
\end{enumerate}
$m$-મૂલ્ય ગોડેલ તર્કશાસ્ત્રની વ્યાખ્યા પણ આવી જ છે;
માત્ર $V = V_m$ છે.
\end{defn}
\sourcecorrection{OLMV-015}{સ્થિર અંગ્રેજી મૂળના
નિષેધ-વિધેયની \texttt{cases} રચના પહેલેથી ગણિતીય
વાતાવરણમાં હોવા છતાં પરિણામો આસપાસ ફરી
\texttt{\$} ચિહ્ન મૂકે છે. તે અયોગ્ય ગણિતીય
વાતાવરણ દૂર કરીને $1$ અને $0$નાં મૂલ્યો
યથાવત રાખ્યાં છે.}

\begin{prop}
  \olref[thr][god]{defn:goedel} વડે વ્યાખ્યાયિત
  $\LogGod[3]$ એ \olref{def:goedel} વડે વ્યાખ્યાયિત
  $\LogGod[3]$ જેવું જ છે.
\end{prop}

\begin{proof}
  \olref[thr][god]{defn:goedel}માં આપેલી સંયોજકોની
  સત્યતાસારણીઓની સરખામણી \olref{def:goedel}નાં
  સમીકરણોથી મળતી સત્યતાસારણીઓ સાથે કરીને આ જોઈ શકાય છે:
  \begin{center}
    \begin{tabular}{c|c} 
      $\tf{\lnot}[\LogGod[3]]$ & \\ 
      \hline  
      $1$ & $0$ \\ 
      $1/2$ & $0$ \\
      $0$ & $1$ 
    \end{tabular}
    \quad
    \begin{tabular}{c|ccc} 
      $\tf{\land}[\LogGod]$ & $1$ & $1/2$ & $0$ \\ 
      \hline 
      $1$ & $1$ & $1/2$ & $0$ \\ 
      $1/2$ & $1/2$ & $1/2$ & $0$\\ 
      $0$ & $0$ & $0$ & $0$ 
    \end{tabular}
    \\[2ex]
    \begin{tabular}{c|ccc} 
      $\tf{\lor}[\LogGod]$ & $1$ & $1/2$ & $0$ \\ 
      \hline 
      $1$ & $1$ & $1$ & $1$ \\ 
      $1/2$ & $1$ & $1/2$ & $1/2$ \\
      $0$ & $1$ & $1/2$ & $0$ 
    \end{tabular}
    \quad
    \begin{tabular}{c|ccc} 
      $\tf{\lif}[\LogGod]$ & $1$ & $1/2$ & $0$ \\ 
      \hline 
      $1$ & $1$ & $1/2$ & $0$ \\ 
      $1/2$ & $1$ & $1$ & $0$  \\ 
      $0$ & $1$ & $1$ & $1$ 
    \end{tabular}
  \end{center} 
\end{proof}

\begin{prop}\ollabel{prop:god-infty-m}
  જો $\Gamma \Entails[\LogGod[\infty]] !B$, તો બધાં~$m \ge 2$
  માટે $\Gamma \Entails[\LogGod[m]] !B$.
\end{prop}

\begin{proof}
  અભ્યાસપ્રશ્ન.
\end{proof}

\begin{prob}
  \olref[mvl][inf][god]{prop:god-infty-m} સાબિત કરો.
\end{prob}

હકીકતમાં, સાન્ત આધારગણ માટે વિપરીત દિશા પણ સાચી છે.
\sourcecorrection{OLMV-016}{સ્થિર અંગ્રેજી મૂળ
વિપરીત દિશા કોઈ મર્યાદા વિના કહે છે. સાન્ત
આધારગણ માટે સંમેય વિરોધી મૂલ્યાંકનમાં આવતા
મર્યાદિત સત્યમૂલ્યો એક જ સાન્ત જાળીમાં સમાય છે.
પરંતુ અનંત આધારગણ ગોડેલના પ્રેરણ વડે અસીમ
ઉતરતા મૂલ્યો ફરજિયાત કરી શકે; તે દરેક સાન્ત
જાળીમાં નિષ્કર્ષને નિર્દિષ્ટ કરે છતાં અનંત
સંમેય જાળીમાં વિરોધી મૂલ્યાંકન આપે છે.
તેથી અહીં દાવો સાન્ત આધારગણ સુધી મર્યાદિત કર્યો છે.}

$\LogGod[3]$ની જેમ $\LogGod[\infty]$માં બધી સ્ફુરણાવાદી રીતે
પ્રામાણ્ય સૂત્રો પુનરુક્તિઓ છે, અને નીચે દર્શાવેલી
બિનપુનરુક્તિઓ~$\LogGod[\infty]$માં પણ બિનપુનરુક્તિઓ છે:
\begin{align*}
  & p \lor \lnot p && (p \lif q) \lif (\lnot p \lor q) \\
  & \lnot\lnot p \lif p && \lnot(\lnot p \land \lnot q) \lif (p \lor q) \\
  & ((p \lif q) \lif p) \lif p && \lnot(p \lif q) \lif (p \land \lnot q)
\end{align*}
સ્ફુરણાવાદી રીતે અપ્રામાણ્ય સૂત્ર હોવા છતાં $\LogGod[3]$ની પુનરુક્તિ
હોવાનું ઉદાહરણ $(p \lif q) \lor (q \lif p)$ એ
$\LogGod[\infty]$માં પણ પુનરુક્તિ છે. હકીકતમાં,
$\LogGod[\infty]$ને $(!A \lif !B) \lor (!B \lif !A)$ની
યોજના ઉમેરેલા સ્ફુરણાવાદી તર્કશાસ્ત્ર તરીકે લક્ષણિત
કરી શકાય છે. આ માઇકલ ડમેટે બતાવ્યું હતું, અને તેથી
$\LogGod[\infty]$ને ગોડેલ--ડમેટ તર્કશાસ્ત્ર~$\Log{LC}$ પણ કહે છે.

\begin{prob}
  $(p \lif q) \lor (q \lif p)$ એ~$\LogGod[\infty]$ની
  પુનરુક્તિ છે તે બતાવો.
\end{prob}

\begin{prob}
  $(p \lif q) \lor (q \lif r) \lor (r \lif s)$ એ
  $\LogGod[3]$ની પુનરુક્તિ છે, પરંતુ~$\LogGod[\infty]$ની
  પુનરુક્તિ નથી, તે બતાવો.
\end{prob}
% END OLP-0401
\endgroup


\OLEndChapterHook
% END OLP-0398
\endgroup


\begingroup
\makeatletter\def\ol@sectioncs{section}\makeatother

% BEGIN OLP-0402
\olchapter{mvl}{seq}{સિક્વન્ટ કલન}
\olchapterid{mvl}{seq}
\phantomsection\label{guunit:OLP-0402}


\begingroup
\makeatletter\def\ol@sectioncs{section}\makeatother

% BEGIN OLP-0403
\olfileid{mvl}{seq}{int}

\olsection{પરિચય}
\phantomsection\label{guunit:OLP-0403}


શાસ્ત્રીય તર્કશાસ્ત્રનું સિક્વન્ટ કલન કાર્યક્ષમ અને સરળ
નિષ્પત્તિ તંત્ર છે. જો કોઈ બહુમૂલ્ય તર્કશાસ્ત્ર સાન્ત
સત્યમૂલ્યો ધરાવતી શ્રેણિક વડે વ્યાખ્યાયિત થાય, એટલે કે $V$
સાન્ત હોય, તો તેના માટે સિક્વન્ટ કલન આપી શકાય છે.
તે કેવી રીતે કરવું તેનો વિચાર શાસ્ત્રીય કિસ્સામાં સિક્વન્ટોના
અર્થ અને અનુમાનના નિયમોના સ્વરૂપ પરથી આવે છે.

હવે યાદ કરો કે સિક્વન્ટ
\begin{align*}
    !A_1, \dots, !A_m & \Sequent !B_1, \dots, !B_n
\intertext{નું અર્થઘટન નીચેના સૂત્ર તરીકે થઈ શકે છે:}
    (!A_1 \land \cdots \land !A_m) & \lif (!B_1 \lor \cdots \lor
!B_n)
\end{align*}
\sourcecorrection{OLMV-017}{સ્થિર અંગ્રેજી મૂળ
સિક્વન્ટની ડાબી બાજુનો છેલ્લો ક્રમસૂચક $n$ લખે છે,
પણ તરત નીચેના અર્થઘટનમાં ડાબી બાજુ માટે $m$ વાપરે છે.
બંને બાજુની લંબાઈ સ્વતંત્ર રહે તે માટે અહીં ડાબે
$!A_1,\dots,!A_m$ અને જમણે $!B_1,\dots,!B_n$ રાખ્યા છે.}
બીજા શબ્દોમાં, સત્યમૂલ્ય-નિયુક્તિ~$\pAssign{v}$ સિક્વન્ટ
$\Gamma \Sequent \Delta$ને \emph{સંતોષે છે} જો અને માત્ર જો
કોઈ $!A \in \Gamma$ માટે $\pValue{v}(!A) = \False$ હોય
અથવા કોઈ $!A \in \Delta$ માટે $\pValue{v}(!A) = \True$ હોય.
આ અર્થઘટન મુજબ આરંભિક સિક્વન્ટો $!A \Sequent !A$
હંમેશા સંતોષાય છે, કારણ કે $\pValue v(!A) = \True$ અથવા
$\pValue{v}(!A) = \False$ હોય છે.
\sourcecorrection{OLMV-018}{સ્થિર અંગ્રેજી મૂળ
આરંભિક સિક્વન્ટ અંગેના અસત્ય-કિસ્સામાં
\texttt{\textbackslash pValue}નું મૂલ્યાંકન-ચલ $v$
છોડી દે છે. બંને વિકલ્પોમાં એ જ $v$ રાખ્યું છે.}

અહીં બાજુનાં સૂત્રો $\Gamma$ અને $\Delta$ છોડી દઈને
$\Log{LK}$માં પ્રેરણ સંયોજકના અનુમાન-નિયમો આપ્યાં છે:

\begin{defish}
    \Axiom$ \fCenter !A$
    \Axiom$ !B \fCenter $
    \RightLabel{\LeftR{\lif}}
    \BinaryInf$ !A \lif !B \fCenter $
    \DisplayProof
    \hfill
    \Axiom$ !A \fCenter !B$
    \RightLabel{\RightR{\lif}}
    \UnaryInf$ \fCenter !A \lif !B $
    \DisplayProof
\end{defish}

ઉપર આપેલા સિક્વન્ટના અર્થઘટનને લાગુ કરીએ તો
$\LeftR{\lif}$ નિયમ કહે છે કે જો $\pValue v(!A) = \True$
અને $\pValue v(!B) = \False$ હોય, તો
$\pValue v(!A \lif !B) = \False$ હોય છે. તેવી જ રીતે,
$\RightR{\lif}$ નિયમ કહે છે કે જો $\pValue v(!A) = \False$
અથવા $\pValue v(!B) = \True$ હોય, તો
$\pValue v(!A \lif !B) = \True$ હોય છે. હકીકતમાં,
આ બંને શરતી વિધાનો દ્વિશરતી છે. $\LeftR{\land}$ અને
$\RightR{\lor}$ નિયમોના પ્રમાણભૂત સ્વરૂપમાં તેમના
અનુરૂપ શરતી વિધાનો દ્વિશરતી નહીં હોય. પરંતુ આ
નિયમોનાં વૈકલ્પિક સ્વરૂપો છે જેમાં તે દ્વિશરતી બને છે:

\begin{defish}
    \Axiom$!A, !B, \Gamma \fCenter \Delta$
    \RightLabel{\LeftR{\land}}
    \UnaryInf$!A \land !B, \Gamma \fCenter \Delta$
    \DisplayProof
    \hfill
    \Axiom$ \Gamma \fCenter \Delta, !A, !B$
    \RightLabel{\RightR{\lor}}
    \UnaryInf$ \Gamma \fCenter \Delta, !A \lor !B$
    \DisplayProof
\end{defish}

આ મૂળભૂત વિચાર $n$-મૂલ્ય તર્કશાસ્ત્ર પર લાગુ કરતાં
બેને બદલે $n$ સ્થાનોવાળું સિક્વન્ટ કલન મળે છે, જેમાં
દરેક સત્યમૂલ્ય માટે એક સ્થાન હોય છે. $V = \{\False, \Undef, \True\}$
ધરાવતા ત્રિમૂલ્ય તર્કશાસ્ત્ર માટે સિક્વન્ટ
$\Gamma \mid \Pi \mid \Delta$ સ્વરૂપનું પદ છે. તે
સત્યમૂલ્ય-નિયુક્તિ~$\pAssign v$માં સંતોષાય છે જો અને માત્ર જો
કોઈ $!A \in \Gamma$ માટે $\pValue{v}(!A) = \False$
અથવા કોઈ $!A \in \Delta$ માટે $\pValue{v}(!A) = \True$
અથવા કોઈ $!A \in \Pi$ માટે $\pValue{v}(!A) = \Undef$
હોય. તેથી આરંભિક સિક્વન્ટો $!A \mid !A \mid !A$
હંમેશા સંતોષાય છે.
% END OLP-0403
\endgroup


\begingroup
\makeatletter\def\ol@sectioncs{section}\makeatother

% BEGIN OLP-0404
\olfileid{mvl}{seq}{rul}

\olsection{નિયમો અને નિષ્પત્તિઓ}
\phantomsection\label{guunit:OLP-0404}


નીચેના માટે, $\Gamma, \Delta, \Pi, \Lambda$ એ વાક્યોની
સાન્ત શ્રેણીઓ દર્શાવે તેમ માનો.

\begin{defn}[સિક્વન્ટ]
\emph{$n$-બાજુવાળો સિક્વન્ટ} નીચેના સ્વરૂપનું પદ છે:
\[
\Gamma_1 \nSequent \dots \nSequent \Gamma_n
\]
અહીં દરેક $\Gamma_i$ એ ભાષા~$\Lang L$નાં વાક્યોની
સાન્ત (સંભવતઃ ખાલી) શ્રેણી છે.
\end{defn}
\sourcecorrection{OLMV-019}{સ્થિર અંગ્રેજી મૂળ
સિક્વન્ટનાં $n$ સ્થાનો દર્શાવ્યા પછી દરેક સ્થાનને
$\Gamma_1$ કહે છે. સામાન્ય સ્થાન માટે $\Gamma_i$
વાપરીને તેની વ્યાખ્યા અને ગણિતીય સૂચકાંક એકરૂપ કર્યા છે.}

\begin{defn}[આરંભિક સિક્વન્ટ]
\emph{$n$-બાજુવાળો આરંભિક સિક્વન્ટ} ભાષાના કોઈપણ
વાક્ય $!A$ માટે $!A \nSequent \dots \nSequent !A$
સ્વરૂપનો $n$-બાજુવાળો સિક્વન્ટ છે.

જો ભાષામાં $0$-સ્થાનવાળો સંયોજક~$\star$ હોય, એટલે
કે વિધાનાત્મક અચલ હોય, તો $\dots \nSequent
\star \nSequent \dots$ સિક્વન્ટ પણ આરંભિક ગણીએ છીએ.
તેમાં $\star$ તેના $\tf{\star} \in V$ સાથે સંકળાયેલા
સત્યમૂલ્યના સ્થાનમાં આવે છે અને અન્ય સ્થાનો ખાલી રહે છે.
\end{defn}

$n$-મૂલ્ય તર્કશાસ્ત્ર~$\Log{L}$ના દરેક સંયોજક માટે,
તે સંયોજક~$\Log{L}$માં લઈ શકે એવા દરેક સત્યમૂલ્યને
અનુરૂપ એક તાર્કિક નિયમ હોય છે. $\Log{L}$ના $n$-બાજુવાળા
સિક્વન્ટ કલનમાં નિષ્પત્તિઓ એ સિક્વન્ટોનાં વૃક્ષો છે:
સૌથી ઉપરના સિક્વન્ટો આરંભિક સિક્વન્ટો હોય છે, અને
જો કોઈ સિક્વન્ટ બીજા એક કે વધુ સિક્વન્ટોની નીચે હોય,
તો તે $\Log{L}$ના સંયોજકોના અનુમાન-નિયમ વડે તેમાંથી
યોગ્ય રીતે મળવું જોઈએ.

\begin{defn}[પ્રમેયો]
વાક્ય~$!A$ એ $n$-મૂલ્ય તર્કશાસ્ત્ર~$\Log{L}$નું
\emph{પ્રમેય} છે જો $\Log{L}$ના નિર્દિષ્ટ સત્યમૂલ્યને
અનુરૂપ દરેક સ્થાનમાં $!A$ ધરાવતા $n$-સિક્વન્ટની
નિષ્પત્તિ હોય. જો $!A$ પ્રમેય હોય તો
$\Proves[\Log{L}] !A$ લખીએ છીએ, અને ન હોય તો
$\Proves/[\Log{L}] !A$ લખીએ છીએ.
\end{defn}

\begin{defn}[નિષ્પન્નક્ષમતા]
વાક્ય~$!A$ એ વાક્યોના ગણ~$\Gamma$માંથી
$n$-મૂલ્ય તર્કશાસ્ત્ર~$\Log{L}$માં \emph{નિષ્પન્ન કરી શકાય એવું}
છે, $\Gamma \Proves[\Log{L}] !A$, જો અને માત્ર જો
કોઈ સાન્ત ઉપગણ~$\Gamma_0 \subseteq \Gamma$
અને $\Gamma_0$નાં વાક્યોની કોઈ શ્રેણી $\Gamma_0'$
હોય જેથી નીચેના સિક્વન્ટની નિષ્પત્તિ હોય:
\[ \Lambda_1 \nSequent \dots \nSequent \Lambda_n \]
અહીં જો સ્થાન $i$ કોઈ નિર્દિષ્ટ સત્યમૂલ્યને અનુરૂપ હોય
તો $\Lambda_i$ એ $!A$ છે, અને અન્યથા $\Gamma_0'$ છે.
જો $!A$ એ~$\Gamma$માંથી નિષ્પન્ન કરી શકાય એવું ન હોય તો
$\Gamma \Proves/ !A$ લખીએ છીએ.
\end{defn}

દાખલા તરીકે, $3$-મૂલ્ય લુકાશેવિચ તર્કશાસ્ત્રમાં
$3$-બાજુવાળું સિક્વન્ટ કલન છે. $3$-બાજુવાળા સિક્વન્ટ
$\Gamma \nSequent \Pi \nSequent \Delta$માં $\Gamma$
એ~$\False$ને, $\Delta$ એ~$\True$ને, અને $\Pi$
એ~$\Undef$ને અનુરૂપ છે. સ્વયંસિદ્ધિઓ
$!A \nSequent !A \nSequent !A$ છે. માત્ર $\True$
નિર્દિષ્ટ હોવાથી, $\Gamma \Proves[\LogLuk[3]] !A$
ત્યારે અને માત્ર ત્યારે જ જ્યારે
$\Gamma \nSequent \Gamma \nSequent !A$ સિક્વન્ટની
નિષ્પત્તિ હોય. (જો $\Undef$ પણ નિર્દિષ્ટ હોત,
તો $\Gamma \nSequent !A \nSequent !A$ની
નિષ્પત્તિ જોઈતી હોત.)
% END OLP-0404
\endgroup


\begingroup
\makeatletter\def\ol@sectioncs{section}\makeatother

% BEGIN OLP-0405
\olfileid{mvl}{seq}{str}

\olsection{સંરચનાત્મક નિયમો}
\phantomsection\label{guunit:OLP-0405}


$n$-બાજુવાળા સિક્વન્ટ કલનના સંરચનાત્મક નિયમો
શાસ્ત્રીય કિસ્સાની જેમ જ કાર્ય કરે છે, પરંતુ દરેક
સ્થાન~$i$ માટે અલગ હોય છે.

\begin{defish}
\begin{center}
\AxiomC{$ \Gamma_1 \nSequent \dots \nSequent \phantom{!A,}\Gamma_i \nSequent \dots \nSequent \Gamma_n $}
\RightLabel{$\iR{\Weakening}{i}$}
\UnaryInfC{$\Gamma_1 \nSequent \dots \nSequent !A, \Gamma_i \nSequent \dots \nSequent \Gamma_n$}
\DisplayProof
\\[2ex]
\AxiomC{$ \Gamma_1 \nSequent \dots \nSequent !A, !A, \Gamma_i \nSequent \dots \nSequent \Gamma_n $}
\RightLabel{$\iR{\Contraction}{i}$}
\UnaryInfC{$\Gamma_1 \nSequent \dots \nSequent \phantom{!A,}!A, \Gamma_i \nSequent
\dots \nSequent \Gamma_n$}
\DisplayProof
\\[2ex]
\AxiomC{$ \Gamma_1 \nSequent \dots \nSequent \Gamma_i, !A, !B, \Gamma_i' \nSequent \dots \nSequent \Gamma_n $}
\RightLabel{$\iR{\Exchange}{i}$}
\UnaryInfC{$\Gamma_1 \nSequent \dots \nSequent \Gamma_i, !B, !A, \Gamma_i' \nSequent \dots
\nSequent \Gamma_n$}
\DisplayProof
\end{center}
\end{defish}

શિથિલીકરણ, સંકોચન અને અદલાબદલીનાં અનુમાનોની શ્રેણી
ઘણી વાર બેવડી અનુમાન-રેખાઓથી દર્શાવવામાં આવશે.

\Cut{} નિયમનાં અનેક સ્વરૂપો છે: સિક્વન્ટમાંનાં ભિન્ન
સ્થાનોના દરેક સંયોજન $i \neq j$ માટે એક સ્વરૂપ.
\begin{defish}
\[
\AxiomC{$ \Gamma_1 \nSequent \dots \nSequent !A, \Gamma_i \nSequent \dots \nSequent \Gamma_n $}
\AxiomC{$ \Delta_1 \nSequent \dots \nSequent !A, \Delta_j \nSequent \dots \nSequent \Delta_n $}
\RightLabel{$\iR{\Cut}{i,j}$}
\BinaryInfC{$\Gamma_1,\Delta_1 \nSequent \dots \nSequent \Gamma_n, \Delta_n$}
\DisplayProof
\]
\end{defish}
% END OLP-0405
\endgroup


\begingroup
\makeatletter\def\ol@sectioncs{section}\makeatother

% BEGIN OLP-0406
\olfileid{mvl}{seq}{prl}

\olsection{પસંદ કરેલાં તર્કશાસ્ત્રો માટે વિધાનાત્મક નિયમો}
\phantomsection\label{guunit:OLP-0406}


$n$-બાજુવાળા સિક્વન્ટ કલનમાં કોઈ સંયોજકના અનુમાન-નિયમો
માત્ર તે સંયોજકના લાક્ષણિક સત્યવિધેય પર આધાર રાખે છે.
તેથી જુદા તર્કશાસ્ત્રોમાં કોઈ સંયોજક એ જ સત્યવિધેય વડે
વ્યાખ્યાયિત હોય, તો તે સંયોજકના $n$-બાજુવાળા સિક્વન્ટ
નિયમો પણ તે તર્કશાસ્ત્રોમાં સમાન હોય છે.

\subsection{$\lnot$ માટેના નિયમો}

નીચેના $\lnot$ માટેના નિયમો લુકાશેવિચ અને ક્લીનીનાં
તર્કશાસ્ત્રો તથા તેમનાં ભેદોને લાગુ પડે છે.

\begin{defish}
  \begin{center}
\AxiomC{$ \Gamma \nSequent \Pi \nSequent \Delta, !A $}
\RightLabel{\iR{\lnot}{\False}}
\UnaryInfC{$ \lnot !A, \Gamma \nSequent \Pi \nSequent \Delta$}
\DisplayProof
\\[2ex]
\AxiomC{$ \Gamma \nSequent !A, \Pi \nSequent \Delta $}
\RightLabel{\iR{\lnot}{\Undef}}
\UnaryInfC{$\Gamma \nSequent \lnot !A, \Pi \nSequent \Delta$}
\DisplayProof
\\[2ex]
\AxiomC{$!A, \Gamma \nSequent \Pi \nSequent \Delta$}
\RightLabel{\iR{\lnot}{\True}}
\UnaryInfC{$\Gamma \nSequent \Pi \nSequent \Delta,  \lnot !A$}
\DisplayProof
  \end{center}
\end{defish}

નીચેના $\lnot$ માટેના નિયમો ગોડેલ તર્કશાસ્ત્રને લાગુ પડે છે.

\begin{defish}
\AxiomC{$ \Gamma \nSequent !A, \Pi \nSequent \Delta, !A $}
\RightLabel{\iR{\lnot}{\False}[\LogGod]}
\UnaryInfC{$ \lnot !A, \Gamma \nSequent \Pi \nSequent \Delta$}
\DisplayProof
\hfill
\AxiomC{$!A, \Gamma \nSequent \Pi \nSequent \Delta$}
\RightLabel{\iR{\lnot}{\True}[\LogGod]}
\UnaryInfC{$\Gamma \nSequent \Pi \nSequent \Delta,  \lnot !A$}
\DisplayProof
\hfill
\end{defish}

(ગોડેલ તર્કશાસ્ત્રમાં $\lnot !A$નું મૂલ્ય ક્યારેય~$\Undef$
હોઈ શકતું નથી; તેથી મધ્યના સ્થાન માટે નિયમ નથી.)

\subsection{$\land$ માટેના નિયમો}

આ $\land$ માટેના નિયમો લુકાશેવિચ, પ્રબળ ક્લીની અને
ગોડેલ તર્કશાસ્ત્રમાં લાગુ પડે છે.

\begin{defish}
\begin{center}
  \AxiomC{$!A, !B, \Gamma \nSequent \Pi \nSequent \Delta$}
  \RightLabel{$\iR{\land}{\False}$}
  \UnaryInfC{$!A \land !B, \Gamma \nSequent \Pi \nSequent \Delta$}
  \DisplayProof
\\[2ex]
  \AxiomC{$\Gamma \nSequent !A, \Pi \nSequent !A, \Delta$}
  \AxiomC{$\Gamma \nSequent !B, \Pi \nSequent !B, \Delta$}
  \AxiomC{$\Gamma \nSequent !A, !B, \Pi \nSequent \Delta$}
  \RightLabel{$\iR\land\Undef$}
  \TrinaryInfC{$\Gamma \nSequent !A \land !B, \Pi \nSequent \Delta$}
  \DisplayProof
  \\[2ex]  
  \AxiomC{$\Gamma \nSequent \Pi \nSequent \Delta, !A$}
  \AxiomC{$\Gamma \nSequent \Pi \nSequent \Delta, !B$}
  \RightLabel{$\iR\land\True$}
  \BinaryInfC{$\Gamma \nSequent \Pi \nSequent \Delta, !A \land !B$}
  \DisplayProof
\end{center}
\end{defish}

\subsection{$\lor$ માટેના નિયમો}

આ $\lor$ માટેના નિયમો લુકાશેવિચ, પ્રબળ ક્લીની અને
ગોડેલ તર્કશાસ્ત્રમાં લાગુ પડે છે.

\begin{defish}
\begin{center}
 \AxiomC{$!A, \Gamma \nSequent \Pi \nSequent \Delta$}
  \AxiomC{$!B, \Gamma \nSequent \Pi \nSequent \Delta$}
  \RightLabel{$\iR\lor\False$}
  \BinaryInfC{$!A \lor !B, \Gamma \nSequent \Pi \nSequent \Delta$}
  \DisplayProof
\\[2ex]
  \AxiomC{$!A, \Gamma \nSequent !A, \Pi \nSequent \Delta$}
  \AxiomC{$!B, \Gamma \nSequent !B, \Pi \nSequent \Delta$}
  \AxiomC{$\Gamma \nSequent !A, !B, \Pi \nSequent \Delta$}
  \RightLabel{$\iR\lor\Undef$}
  \TrinaryInfC{$\Gamma \nSequent !A \lor !B, \Pi \nSequent \Delta$}
  \DisplayProof
  \\[2ex]  
  \AxiomC{$\Gamma \nSequent \Pi \nSequent \Delta, !A, !B$}
  \RightLabel{$\iR\lor\True$}
  \UnaryInfC{$\Gamma \nSequent \Pi \nSequent \Delta, !A \lor !B$}
  \DisplayProof
\end{center}
\end{defish}

\subsection{$\lif$ માટેના નિયમો}

આ $\lif$ માટેના નિયમો લુકાશેવિચ તર્કશાસ્ત્રમાં લાગુ પડે છે.

\begin{defish}
\begin{center}
  \AxiomC{$\Gamma \nSequent \Pi \nSequent \Delta, !A$}
  \AxiomC{$!B, \Gamma \nSequent \Pi \nSequent \Delta$}
  \RightLabel{$\iR\lif\False[\LogLuk[3]]$}
  \BinaryInfC{$!A \lif !B, \Gamma \nSequent \Pi \nSequent \Delta$}
  \DisplayProof
  \\[2ex]
  \AxiomC{$\Gamma \nSequent !A, !B, \Pi \nSequent \Delta$}
  \AxiomC{$!B, \Gamma \nSequent \Pi \nSequent \Delta, !A$}
  \RightLabel{$\iR\lif\Undef[\LogLuk[3]]$}
  \BinaryInfC{$\Gamma \nSequent !A \lif !B, \Pi \nSequent \Delta$}
  \DisplayProof
\\[2ex]
  \AxiomC{$!A, \Gamma \nSequent !B, \Pi \nSequent \Delta, !B$}
  \AxiomC{$!A, \Gamma \nSequent !A, \Pi \nSequent \Delta, !B$}
  \RightLabel{$\iR{\lif}{\True}[\LogLuk[3]]$}
  \BinaryInfC{$\Gamma \nSequent \Pi \nSequent \Delta, !A \lif !B$}
  \DisplayProof
\end{center}
\end{defish}

આ $\lif$ માટેના નિયમો પ્રબળ ક્લીની તર્કશાસ્ત્રમાં લાગુ પડે છે.

\begin{defish}
\begin{center}
  \AxiomC{$\Gamma \nSequent \Pi \nSequent \Delta, !A$}
  \AxiomC{$!B, \Gamma \nSequent \Pi \nSequent \Delta$}
  \RightLabel{$\iR\lif\False[\LogKs]$}
  \BinaryInfC{$!A \lif !B, \Gamma \nSequent \Pi \nSequent \Delta$}
  \DisplayProof
  \\[2ex]
  \AxiomC{$!B, \Gamma \nSequent !B, \Pi \nSequent \Delta$}
  \AxiomC{$\Gamma \nSequent !A, !B, \Pi \nSequent \Delta$}
  \AxiomC{$\Gamma \nSequent !A, \Pi \nSequent \Delta, !A$}
  \RightLabel{$\iR\lif\Undef[\LogKs]$}
  \TrinaryInfC{$\Gamma \nSequent !A \lif !B, \Pi \nSequent \Delta$}
  \DisplayProof
\\[2ex]
  \AxiomC{$!A, \Gamma \nSequent \Pi \nSequent \Delta, !B$}
  \RightLabel{$\iR{\lif}{\True}[\LogKs]$}
  \UnaryInfC{$\Gamma \nSequent \Pi \nSequent \Delta, !A \lif !B$}
  \DisplayProof
\end{center}
\end{defish}

આ $\lif$ માટેના નિયમો ગોડેલ તર્કશાસ્ત્રમાં લાગુ પડે છે.

\begin{defish}
\begin{center}
  \AxiomC{$\Gamma \nSequent !A, \Pi \nSequent \Delta, !A$}
  \AxiomC{$!B, \Gamma \nSequent \Pi \nSequent \Delta$}
  \RightLabel{$\iR\lif\False[\LogGod[3]]$}
  \BinaryInfC{$!A \lif !B, \Gamma \nSequent \Pi \nSequent \Delta$}
  \DisplayProof
  \\[2ex]
  \AxiomC{$\Gamma \nSequent !B, \Pi \nSequent \Delta$}
  \AxiomC{$\Gamma \nSequent \Pi \nSequent \Delta, !A$}
  \RightLabel{$\iR\lif\Undef[\LogGod[3]]$}
  \BinaryInfC{$\Gamma \nSequent !A \lif !B, \Pi \nSequent \Delta$}
  \DisplayProof
\\[2ex]
  \AxiomC{$!A, \Gamma \nSequent !B, \Pi \nSequent \Delta, !B$}
  \AxiomC{$!A, \Gamma \nSequent !A, \Pi \nSequent \Delta, !B$}
  \RightLabel{$\iR{\lif}{\True}[\LogGod[3]]$}
  \BinaryInfC{$\Gamma \nSequent \Pi \nSequent \Delta, !A \lif !B$}
  \DisplayProof
\end{center}
\end{defish}


\begin{sidewaysfigure}
  \begin{center}  
  \AxiomC{$A \nSequent A \nSequent A$}
  \RightLabel{\iR \Weakening \True}
  \UnaryInfC{$A \nSequent A \nSequent B, A$}
  \RightLabel{\iR \Weakening \Undef}
  \UnaryInfC{$A \nSequent B, A \nSequent B, A$}
  \RightLabel{\iR \Weakening \Undef}
  \UnaryInfC{$A \nSequent A, B, A \nSequent B, A$}
  \AxiomC{$A \nSequent A \nSequent A$}
  \RightLabel{\iR \Weakening \True}
  \UnaryInfC{$A \nSequent A \nSequent A, A$}
  \RightLabel{\iR \Weakening \True}
  \UnaryInfC{$A \nSequent A \nSequent B, A, A$}
  \RightLabel{\iR \Weakening \False}
  \UnaryInfC{$B, A \nSequent A \nSequent B, A, A$}
  \RightLabel{$\iR\lif\Undef$}
  \BinaryInfC{$A \nSequent A \lif B, A \nSequent B, A$}
  \AxiomC{$B \nSequent B \nSequent B$}
  \RightLabel{\iR \Weakening \Undef}
  \UnaryInfC{$B \nSequent A, B \nSequent B$}
  \RightLabel{\iR \Exchange \Undef}
  \UnaryInfC{$B \nSequent B, A \nSequent B$}
  \RightLabel{\iR \Weakening \Undef}
  \UnaryInfC{$B \nSequent A, B, A \nSequent B$}
  \RightLabel{\iR \Weakening \False}
  \UnaryInfC{$A, B \nSequent A, B, A \nSequent B$}
  \RightLabel{\iR \Exchange \False}
  \UnaryInfC{$B, A \nSequent A, B, A \nSequent B$}
  \AxiomC{$A \nSequent A \nSequent A$}
  \RightLabel{\iR \Weakening \True}
  \UnaryInfC{$A \nSequent A \nSequent B, A$}
  \RightLabel{\iR \Weakening \False}
  \UnaryInfC{$B, A \nSequent A \nSequent B, A$}
  \RightLabel{\iR \Weakening \False}
  \UnaryInfC{$B, B, A \nSequent A \nSequent B, A$}
  \RightLabel{$\iR\lif\Undef$}
  \BinaryInfC{$B, A \nSequent A \lif B, A \nSequent B$}
  \RightLabel{$\iR\lif\False$}
  \BinaryInfC{$A \lif B, A \nSequent A \lif B, A \nSequent B$}
  \DisplayProof
  \end{center}
  \caption{$\LogLuk[3]$માં નિષ્પત્તિનું ઉદાહરણ}
\end{sidewaysfigure}
% END OLP-0406
\endgroup


\OLEndChapterHook


\begin{editorial}
આગળનો સંબંધિત અંશ મૂળના સક્રિય આયાતક્રમમાં નથી. તેને અહીં સ્પષ્ટ વૈકલ્પિક સંદર્ભમાં જાળવ્યો છે.
\end{editorial}
\begingroup
\makeatletter\def\ol@sectioncs{section}\makeatother

% BEGIN OLP-0652
\olfileid{mvl}{seq}{prf}

\olsection{સાબિતી-સૈદ્ધાંતિક ખ્યાલો}
\phantomsection\label{guunit:OLP-0652}


\begin{explain}
આપણે અનેક મહત્વના અર્થાનુસારી ખ્યાલો
(પ્રામાણ્ય, ફલિતતા અને સંતોષ્યતા) વ્યાખ્યાયિત
કર્યા છે; એ જ રીતે હવે તેમને અનુરૂપ
\emph{સાબિતી-સૈદ્ધાંતિક ખ્યાલો}
વ્યાખ્યાયિત કરીએ છીએ. આ ખ્યાલો
સંરચનાઓમાં વાક્યોના
સંતોષનો આશરો લઈને નહીં, પરંતુ અમુક
સિક્વન્ટોની નિષ્પન્નક્ષમતા અથવા
અનિષ્પન્નક્ષમતાનો આશરો લઈને
વ્યાખ્યાયિત થાય છે. આ બંને પ્રકારના
ખ્યાલો એકરૂપ થાય છે એ એક મહત્વની
શોધ હતી. તેઓ એકરૂપ થાય છે તે જ
\emph{યથાર્થતા} અને \emph{પૂર્ણતાના પ્રમેય}
દર્શાવે છે.
\end{explain}

\begin{defn}[પ્રમેયો]
જો સિક્વન્ટ $\quad \Sequent !A$ની
$\Log{LK}$માં કોઈ નિષ્પત્તિ હોય,
તો વાક્ય~$!A$ \emph{પ્રમેય}
છે. $!A$ પ્રમેય હોય તો $\Proves
!A$ અને ન હોય તો $\Proves/ !A$
લખીએ છીએ.
\end{defn}

\begin{defn}[નિષ્પન્નક્ષમતા]
વાક્ય~$!A$ એ
વાક્યોના ગણ~$\Gamma$માંથી
\emph{નિષ્પન્ન કરી શકાય એવું} છે, એટલે
$\Gamma \Proves !A$, જો અને માત્ર જો
એવો સાન્ત ઉપગણ~$\Gamma_0 \subseteq \Gamma$
અને $\Gamma_0$માંનાં વાક્યોની
એવી શ્રેણી $\Gamma_0'$ હોય કે
$\Log{LK}$ એ $\Gamma_0' \Sequent !A$ને
નિષ્પન્ન કરે. $!A$ એ $\Gamma$માંથી
નિષ્પન્ન કરી શકાય એવું ન હોય તો
$\Gamma \Proves/ !A$ લખીએ છીએ.
\end{defn}

સંકોચન, શિથિલીકરણ અને અદલાબદલીના
નિયમોને કારણે $\Gamma_0'$માંનાં
વાક્યોનો ક્રમ અને તેમની સંખ્યા
મહત્ત્વ ધરાવતાં નથી: જો
$\Gamma_0' \Sequent !A$ નિષ્પન્ન કરી શકાય એવું
હોય, તો $\Gamma_0'$માં હોય એ જ
વાક્યો ધરાવતી કોઈપણ
$\Gamma_0''$ માટે
$\Gamma_0'' \Sequent !A$ પણ
નિષ્પન્ન થઈ શકે છે. દાખલા તરીકે,
જો $\Gamma_0 = \{!B, !C\}$ હોય,
તો $\Gamma_0' = \tuple{!B, !B, !C}$
અને $\Gamma_0'' =
\tuple{!C, !C, !B}$ બંને
એવી શ્રેણીઓ છે જેમાં માત્ર
$\Gamma_0$નાં વાક્યો જ છે.
તેમાંથી એક ધરાવતું સિક્વન્ટ
નિષ્પન્ન કરી શકાય એવું હોય, તો બીજું
ધરાવતું સિક્વન્ટ પણ નિષ્પન્ન
થઈ શકે છે; દાખલા તરીકે:
\begin{prooftree}
  \AxiomC{}
  \Deduce$!B, !B, !C \fCenter !A$
  \RightLabel{\LeftR{\Contraction}}
  \UnaryInf$!B, !C \fCenter !A$
  \RightLabel{\LeftR{\Exchange}}
  \UnaryInf$!C, !B \fCenter !A$
  \RightLabel{\LeftR{\Weakening}}
  \UnaryInf$!C, !C, !B \fCenter !A$
\end{prooftree}
હવેથી, $\Gamma_0$ એ વાક્યોનો
સાન્ત ગણ હોય ત્યારે
$\Gamma_0 \Sequent !A$ એ એવું
કોઈપણ સિક્વન્ટ છે જેનું પૂર્વાંગ
$\Gamma_0$માંનાં વાક્યોની
શ્રેણી છે, એમ કહીશું; અને જરૂર
પડે ત્યાં સંકોચન, અદલાબદલી
અને શિથિલીકરણોને ગર્ભિત રીતે
સમાવિષ્ટ માનીશું.

\begin{defn}[સુસંગતતા]
વાક્યોનો ગણ~$\Gamma$ \emph{અસુસંગત}
છે જો અને માત્ર જો એવો સાન્ત
ઉપગણ~$\Gamma_0 \subseteq \Gamma$
હોય કે $\Log{LK}$ એ
$\Gamma_0 \Sequent \quad$ને
નિષ્પન્ન કરે. $\Gamma$
અસુસંગત ન હોય, એટલે કે
દરેક સાન્ત $\Gamma_0 \subseteq \Gamma$
માટે $\Log{LK}$ એ
$\Gamma_0 \Sequent \quad$ને
નિષ્પન્ન ન કરે, તો તેને
\emph{સુસંગત} કહીએ છીએ.
\end{defn}

\begin{prop}[સ્વવાચકતા]
\ollabel{prop:reflexivity}
જો $!A \in \Gamma$ હોય, તો
$\Gamma \Proves !A$.
\end{prop}

\begin{proof}
આરંભિક સિક્વન્ટ $!A \Sequent !A$
નિષ્પન્ન કરી શકાય એવું છે અને
$\{!A\} \subseteq \Gamma$ છે.
\end{proof}

\begin{prop}[એકદિશવર્ધિતા]
\ollabel{prop:monotonicity}
જો $\Gamma \subseteq \Delta$ અને
$\Gamma \Proves !A$ હોય, તો
$\Delta \Proves !A$.
\end{prop}

\begin{proof}
માનો કે $\Gamma \Proves !A$,
એટલે એવો સાન્ત $\Gamma_0
\subseteq \Gamma$ છે કે
$\Gamma_0 \Sequent !A$
નિષ્પન્ન કરી શકાય એવું છે.
$\Gamma \subseteq \Delta$
હોવાથી $\Gamma_0$ એ
$\Delta$નો પણ સાન્ત ઉપગણ છે.
તેથી $\Gamma_0 \Sequent !A$ની
નિષ્પત્તિ એ
$\Delta \Proves !A$ પણ બતાવે છે.
\end{proof}

\begin{prop}[પરંપરિતતા]
\ollabel{prop:transitivity}
જો $\Gamma \Proves !A$ અને
$\{!A\} \cup \Delta \Proves
!B$ હોય, તો $\Gamma \cup \Delta
\Proves !B$.
\end{prop}

\begin{proof}
જો $\Gamma \Proves !A$ હોય,
તો કોઈ સાન્ત $\Gamma_0 \subseteq \Gamma$
અને $\Gamma_0 \Sequent !A$ની
નિષ્પત્તિ~$\pi_0$ હોય છે.
જો $\{!A\} \cup \Delta \Proves !B$
હોય, તો કોઈ સાન્ત ઉપગણ
$\Delta_0 \subseteq \Delta$
માટે $!A, \Delta_0 \Sequent !B$ની
નિષ્પત્તિ~$\pi_1$ હોય છે.
નીચેની નિષ્પત્તિ જુઓ:
\begin{prooftree}
  \AxiomC{}
  \RightLabel{$\pi_0$}
  \Deduce$\Gamma_0 \fCenter !A$
  \AxiomC{}
  \RightLabel{$\pi_1$}
  \Deduce$!A, \Delta_0 \fCenter !B$
  \RightLabel{\Cut}
  \BinaryInf$\Gamma_0, \Delta_0 \fCenter !B$
\end{prooftree}
$\Gamma_0 \cup \Delta_0 \subseteq
\Gamma \cup \Delta$ હોવાથી આ
$\Gamma \cup \Delta \Proves !B$
બતાવે છે.
\end{proof}

નોંધો કે ખાસ કરીને તેનો અર્થ એવો
થાય છે કે જો $\Gamma \Proves !A$
અને $!A \Proves !B$ હોય, તો
$\Gamma \Proves !B$. વળી, જો
$!A_1, \dots, !A_n \Proves !B$
અને દરેક~$i$ માટે
$\Gamma \Proves !A_i$ હોય,
તો $\Gamma \Proves !B$
એવું પણ ફલિત થાય છે.

\begin{prop}
\ollabel{prop:incons}
$\Gamma$ અસુસંગત હોય જો અને
માત્ર જો દરેક વાક્ય~$!A$ માટે
$\Gamma \Proves {!A}$.
\end{prop}

\begin{proof}
સ્વાધ્યાય.
\end{proof}

\begin{prob}
\olref[mvl][seq][prf]{prop:incons} સાબિત કરો.
\sourcecorrection{OLMD-189}{મૂળ કસરત
આ વિભાગના \texttt{mvl/seq/prf}
પ્રતિજ્ઞાપનને બદલે
\texttt{fol/seq/ptn}ના સમાન
પ્રતિજ્ઞાપન તરફ નિર્દેશ કરે છે;
કસરતનું સ્થાનિક લક્ષ્ય અહીં રાખ્યું છે.}
\end{prob}

\begin{prop}[સઘનતા]
  \ollabel{prop:proves-compact}
  \begin{enumerate}
  \item જો $\Gamma \Proves !A$ હોય, તો
    એવો સાન્ત ઉપગણ $\Gamma_0
    \subseteq \Gamma$ છે કે
    $\Gamma_0 \Proves !A$.
  \item જો $\Gamma$નો દરેક
    સાન્ત ઉપગણ સુસંગત હોય,
    તો $\Gamma$ સુસંગત છે.
  \end{enumerate}
\end{prop}

\begin{proof}
  \begin{enumerate}
    \item જો $\Gamma \Proves !A$ હોય,
      તો એવો સાન્ત ઉપગણ
      $\Gamma_0 \subseteq \Gamma$ છે કે
      સિક્વન્ટ $\Gamma_0
      \Sequent !A$ને
      નિષ્પત્તિ હોય. તેથી
      $\Gamma_0 \Proves !A$.
    \item જો $\Gamma$ અસુસંગત હોય,
      તો એવો સાન્ત ઉપગણ
      $\Gamma_0 \subseteq \Gamma$ છે કે
      $\Log{LK}$ એ $\Gamma_0
      \Sequent \quad$ને
      નિષ્પન્ન કરે. પરંતુ
      ત્યારે $\Gamma_0$ એ
      $\Gamma$નો અસુસંગત
      સાન્ત ઉપગણ છે.
  \end{enumerate}
\end{proof}
% END OLP-0652
\endgroup
% END OLP-0402
\endgroup


\OLEndPartHook
% END OLP-0383
